%% 4-8 三角関数：円運動の影（第13回：組 V4-5。ページ制限なし。偶数ページ）
%% ほかの巻が指している形：単位円の上の点 (cosθ, sinθ)（5-15）、sin²θ+cos²θ=1（半径1の円で三平方の定理。5-15・5-16）、
%%   90°回すと (x, y)→(−y, x)、cos(θ+90°)=−sinθ、sin(θ+90°)=cosθ（5-15 の「直角三角形をまるごと90°回す」図と同じ見せ方）、
%%   θ=ωt、角速度 ω（rad/s）、周期 T=2π/ω、f=1/T、(r cos ωt, r sin ωt)、v=rω（5-15・物理シリーズ）
\makeatletter
% 小見出し（太字の1行）  \bkivihead{見出し}
\newcommand\bkivihead[1]{\par\addvspace{6pt}\Needspace{6\baselineskip}\noindent{\sffamily\bfseries\color{\bkTT}#1}\par\nobreak\vspace{2pt}}
% 囲みの中の小見出し（それだけで1行）
\newcommand\bkivibox[1]{\par\addvspace{3pt}\noindent\textbf{#1}\par\noindent\ignorespaces}
% 大事な式のわく（巻の色のうすい地に、濃い色の線）  \bkiviframe[はば]{中身}
\newcommand\bkiviframe[2][100mm]{\par\addvspace{4pt}{\centering\tikz\node[draw=\bkTT, line width=0.9pt, fill=bkTheme!10, rounded corners=1.5mm, inner xsep=6mm, inner ysep=2.5mm, align=center]{#2};\par}\addvspace{4pt}\noindent\ignorespaces}
% 例題の途中式の行（折り返した行も、そろえた位置から始まる）
\newcommand\bkiviq{\par\vspace{2pt}\noindent\hangindent=4.3\zw\hangafter=1\hspace*{2.6\zw}}% （1）などで始まる行
\newcommand\bkivia{\par\vspace{2pt}\noindent\hangindent=2.6\zw\hangafter=0\relax}
\newcommand\bkivib{\par\vspace{2pt}\noindent\hangindent=4.3\zw\hangafter=0\relax}
% 式を、まん中に1行で置く
\newcommand\bkiviline[1]{\par\addvspace{3pt}{\centering #1\par}\addvspace{3pt}\noindent\ignorespaces}
% 座標の軸（原点中心、はば #1 mm、高さ #2 mm の半分ずつ）  \bkiviaxes{x の半分}{y の半分}
\newcommand\bkiviaxes[2]{%
  \draw[bkGray, line width=0.6pt, -{Stealth[length=4pt]}] (-#1,0) -- (#1,0) node[anchor=west, font=\sffamily\footnotesize, text=bkInk] {$x$};
  \draw[bkGray, line width=0.6pt, -{Stealth[length=4pt]}] (0,-#2) -- (0,#2) node[anchor=south, font=\sffamily\footnotesize, text=bkInk] {$y$};}
% ゴンドラ  \bkivigondola{(x,y)}{色}
\newcommand\bkivigondola[2]{\begin{scope}[shift={#1}]
  \draw[bkInk!70, line width=0.5pt] (0,0) -- (0,-1.2);
  \filldraw[fill=#2, draw=#2!70!black, line width=0.4pt, rounded corners=0.6mm] (-2,-1.2) rectangle (2,-4.4);
  \fill[white] (-1.3,-1.8) rectangle (1.3,-2.9);
\end{scope}}
\makeatother

\bkchapter{4}{8}{三角関数：\\円運動の影}{\bklevelrange{3}{4}}{三角関数：円運動の影}
\bkmichishirube{1-15,3-9,4-2,4-6,4-7}{5-15,5-16,6-6,7-6}

\begin{goal}
\begin{itemize}[label={\color{\bkTT}●}, leftmargin=1.6\zw, itemsep=0pt, topsep=1pt]
\item \bklevel{3}\enspace 一般角で、単位円の上の点 $(\cos\theta,\ \sin\theta)$ を求められる。
\item \bklevel{3}\enspace 符号・周期、$\sin^2\theta+\cos^2\theta=1$ や90°回す式を、図で説明できる。
\item \bklevel{4}\enspace 角速度 $\omega$ と周期 $T$ で、回る点の位置・速さを求められる。
\end{itemize}
\end{goal}

\begin{hatena}
\noindent\begin{minipage}[c]{0.6\linewidth}
\raggedright
観覧車のゴンドラに真横から光を当てると、うしろのかべに影が映ります。ゴンドラは同じ速さで回るのに、影は、いちばん上と下ではゆっくり、まん中の高さでは速く動きます。

半径\qty{50}{m}、1周20分の観覧車に、いちばん下から乗ります。7分後のゴンドラの高さは何m？　何周もしたあとの105分後は？
\end{minipage}\hfill
\begin{minipage}[c]{0.37\linewidth}
\centering
\begin{tikzpicture}[x=0.6mm, y=0.6mm]
  \draw[bkGray, line width=1pt] (-30,0) -- (46,0);
  % 光（左から真横に進む）
  \foreach \h in {10,24,38} {\draw[bkOrange, line width=0.8pt, -{Stealth[length=4pt]}] (-30,\h) -- (-23,\h);}
  \node[font=\sffamily\footnotesize, anchor=south, text=bkInk] at (-26.5,40.5) {光};
  % 観覧車
  \draw[bkInk!60, line width=0.8pt] (0,24) -- (-9,0) (0,24) -- (9,0);
  \draw[bkInk!60, line width=0.9pt] (0,24) circle (18);
  \foreach \a in {0,45,...,315} {\draw[bkInk!35, line width=0.4pt] (0,24) -- ++(\a:18);}
  \foreach \a in {0,45,90,135,180,225,270,315} {\bkivigondola{($(0,24)+(\a:18)$)}{bkRoute2}}
  \fill[bkInk] (0,24) circle (0.9);
  % 光る1台と、その影
  \bkivigondola{($(0,24)+(45:18)$)}{bkOrange}
  \draw[bkGray, dashed, line width=0.5pt] ($(0,24)+(45:18)+(2.4,-2.8)$) -- (39.4,{24+18*sin(45)-2.8});
  \draw[bkGray!70, line width=2.2pt] (40.5,0) -- (40.5,46);
  \fill[bkOrange!85!black] (39.2,{24+18*sin(45)-4.4}) rectangle (40.5,{24+18*sin(45)-1.2});
  \draw[bkInk, line width=0.7pt, {Stealth[length=3.5pt]}-{Stealth[length=3.5pt]}] (44,{24-18-2.8}) -- (44,{24+18-2.8});
  \node[font=\sffamily\footnotesize, anchor=south, text=bkInk] at (40.5,46.5) {かべ};
\end{tikzpicture}
\end{minipage}
\end{hatena}

\noindent 4-6 と 4-7 では、三角形の角の $\sin$ と $\cos$ を学びました。三角形の角は180°より小さいので、そのままでは、何周も回る観覧車には使えません。この章では、$\sin$ と $\cos$ を「円の上を回る点の位置」として決めなおします。すると、何周回っても、時間がたっても使える、回るもの・ゆれるもの・波の数学の道具になります。

%% ---------------- やってみよう・思い出そう ----------------
\begin{yattemiyou}[回る点の高さをはかろう]
\textbf{用意するもの}\par
方眼紙、コンパス、分度器、定規
\begin{enumerate}[label={\bkcircled[bkGreen]{\arabic*}}, leftmargin=2.4\zw, itemsep=1pt, topsep=3pt]
\item 方眼紙のまん中に点Oをとり、Oを通る横の線とたての線を引こう。Oを中心に、半径\qty{10}{cm}の円をかこう。
\item 分度器で、Oから右へのびる線から左回りに30°ずつ、円の上に点をとろう（0°、30°、……、330°の12こ）。
\item それぞれの点の高さ（横の線から点までの、たての長さ）を定規ではかろう。横の線より下の点は、高さにマイナスをつけて書こう。
\item 角を横軸、高さをたて軸にして点を打ち、なめらかにつなごう。どんな形になるかな？
\end{enumerate}
\begin{bkcenter}
\begin{bktable}{colspec={Q[l,m,wd=22mm] *{6}{Q[c,m,wd=11.5mm]}}, rowsep=1pt, colsep=3pt}
角 & 0° & 30° & 60° & 90° & 120° & 150° \\
高さ（cm） & 0 & 5.0 & 8.7 & 10.0 & 8.7 & 5.0 \\
\end{bktable}\par\vspace{3pt}
\begin{bktable}{colspec={Q[l,m,wd=22mm] *{6}{Q[c,m,wd=11.5mm]}}, rowsep=1pt, colsep=3pt}
角 & 180° & 210° & 240° & 270° & 300° & 330° \\
高さ（cm） & 0 & $-5.0$ & $-8.7$ & $-10.0$ & $-8.7$ & $-5.0$ \\
\end{bktable}
\bktabcaption{ソラさんがはかった結果}
\end{bkcenter}
\end{yattemiyou}

\noindent 30°の点の高さ\qty{5.0}{cm}は、半径\qty{10}{cm}のちょうど半分です。4-6 で学んだ $\sin30°=0.5$ の「0.5」が、ここに出ています。60°の\qty{8.7}{cm}も、$\sin60°$（約0.866）の約10倍です。

④でつないだ形は、山と谷がかわるがわるくる「波」の形になったはずです。180°をこえると高さはマイナスになり、360°で0にもどります。この章では、この波の正体を調べます。

\noindent この章では、「≈」という記号を使います。≈ は「ほぼ等しい」という意味です（5-8 でくわしく学びます）。たとえば $\sqrt3\approx1.732$ は、「$\sqrt3$ はほぼ1.732」と読みます。

\begin{omoidasou}
\bkrecall{1-15}{100日後の曜日：$100\div7=14$ あまり2。7日ごとに同じ曜日にもどるので、あまりの2日だけ考えればよい。}
\bkrecall{4-2}{1周は $360°=2\pi$\,rad。弧の長さ＝半径×角（$s=r\theta$。角はラジアン）。\\電卓は DEG（度）と RAD（ラジアン）を切りかえて使う。\\極座標 $(r,\ \theta)$：原点からの距離 $r$ と、$x$ 軸の正の向きから左回りの角 $\theta$ で場所を言う。}
\bkrecall{4-4}{三平方の定理：直角三角形で（斜辺）$^2$＝（横）$^2$＋（たて）$^2$。\\原点を中心とする半径 $r$ の円の式は $x^2+y^2=r^2$。}
\bkrecall{4-6}{直角三角形で、角 $\theta$ から見て $\sin\theta=\dfrac{\text{たて}}{\text{斜辺}}$、$\cos\theta=\dfrac{\text{横}}{\text{斜辺}}$。\\$\sin30°=0.5$、$\cos30°=\dfrac{\sqrt3}{2}\approx0.866$、$\sin45°=\cos45°=\dfrac{\sqrt2}{2}\approx0.707$。}
\bkrecall{4-6}{極座標 $(r,\ \theta)$ の点を $(x,\ y)$ に直すと、$x=r\cos\theta$、$y=r\sin\theta$。\\半径1の半円の上で、角 $\theta$ の点の座標は $(\cos\theta,\ \sin\theta)$（0°〜180°）。}
\bkrecall{4-7}{鈍角（90°より大きく180°より小さい角）では、$\cos$ は負。}
\bkrecall{3-9}{$y=a(x-p)^2+q$ のグラフは、$y=ax^2$ のグラフを、右へ $p$、上へ $q$ だけ平行移動したもの。}
\end{omoidasou}


%% ---------------- 数学の道具 ----------------
\bksection{doguu}{数学の道具}

\bksubsection{1}{何周しても使える角：一般角}

\leavevmode\label{bk:4-8:ippan}分度器ではかる角は、0°から360°まででした（1-20）。でも、観覧車も、時計の針も、自転車のタイヤも、何周も回りつづけます。回った角をまとめて1つの数で言うには、360°より大きい角が必要です。

たとえば、観覧車が2周して、さらに30°回ったとします。1周は360°なので、回った角は全部で
\bkiviline{$360°\times2+30°=750°$}
です。このときのゴンドラの場所は、30°だけ回った所と同じです。1周回ると、もとの場所にもどるからです。

\begin{bkfigure}
\begin{tikzpicture}[x=1mm, y=1mm]
  % 左：750°（2周と30°）
  \begin{scope}
  \draw[bkGray, line width=0.6pt] (-22,0) -- (22,0);
  \draw[bkGray, line width=0.6pt] (0,-20) -- (0,21);
  \draw[bkBlue, line width=1.1pt, -{Stealth[length=5pt]}, domain=0:750, samples=220, smooth, variable=\t] plot ({(4+\t/72)*cos(\t)},{(4+\t/72)*sin(\t)});
  \draw[bkInk, line width=1.4pt] (0,0) -- (30:19);
  \fill[bkInk] (0,0) circle (0.8);
  \fill[bkOrange] (30:19) circle (1.2);
  \node[font=\sffamily\footnotesize, anchor=west] at (30:20) {P};
  \node[font=\sffamily\footnotesize, align=center, anchor=north] at (0,-22) {750°\\（2周と30°）};
  \end{scope}
  % まん中：30°
  \begin{scope}[xshift=44mm]
  \draw[bkGray, line width=0.6pt] (-22,0) -- (22,0);
  \draw[bkGray, line width=0.6pt] (0,-20) -- (0,21);
  \draw[bkBlue, line width=1.1pt, -{Stealth[length=5pt]}] (8,0) arc (0:30:8);
  \draw[bkInk, line width=1.4pt] (0,0) -- (30:19);
  \fill[bkInk] (0,0) circle (0.8);
  \fill[bkOrange] (30:19) circle (1.2);
  \node[font=\sffamily\footnotesize, anchor=west] at (30:20) {P};
  \node[font=\sffamily\footnotesize] at (15:12.5) {30°};
  \node[font=\sffamily\footnotesize, align=center, anchor=north] at (0,-22) {30°\\（750°と同じ場所）};
  \end{scope}
  % 右：−120°
  \begin{scope}[xshift=88mm]
  \draw[bkGray, line width=0.6pt] (-22,0) -- (22,0);
  \draw[bkGray, line width=0.6pt] (0,-20) -- (0,21);
  \draw[bkRed, line width=1.1pt, -{Stealth[length=5pt]}] (8,0) arc (0:-120:8);
  \draw[bkBlue, line width=0.9pt, dash pattern=on 2.5pt off 1.5pt, -{Stealth[length=4.5pt]}] (13,0) arc (0:240:13);
  \draw[bkInk, line width=1.4pt] (0,0) -- (240:19);
  \fill[bkInk] (0,0) circle (0.8);
  \fill[bkOrange] (240:19) circle (1.2);
  \node[font=\sffamily\footnotesize, anchor=east] at (240:20) {P};
  \node[font=\sffamily\footnotesize, text=bkRed!80!black, anchor=north west, inner sep=1pt] at (6.5,-9) {$-120°$};
  \node[font=\sffamily\footnotesize, text=bkBlue!80!black, anchor=south east, inner sep=1pt] at (150:14.8) {240°};
  \node[font=\sffamily\footnotesize, align=center, anchor=north] at (0,-22) {$-120°$ と 240°\\（同じ場所）};
  \end{scope}
\end{tikzpicture}
\bkfigcaption{左とまん中：750°回った点Pは、30°回った点と同じ場所にいる。右：右回りに120°（$-120°$）回った点は、左回りに240°回った点と同じ場所。}
\end{bkfigure}

これは、1-15で考えた「100日後の曜日」と同じ考え方です。曜日は7日ごとに同じにもどるので、「何週と何日」に分けて、あまりの日だけを見ればよかったのでした。角は360°ごとに同じ場所にもどるので、「何周と何度」に分けて、あまりの角だけを見ればよいのです。

\begin{bkcenter}
\begin{bktable}{colspec={Q[l,m,wd=10mm] Q[c,m,wd=20mm] Q[c,m,wd=48mm] Q[c,m,wd=28mm]}, rowsep=1.5pt, row{1}={font=\sffamily\small}}
 & 1回りのはば & 分け方 & 見ればよいもの \\
曜日 & 7日 & $100\div7=14$ あまり2 & あまりの2日 \\
角 & 360° & $750\div360=2$ あまり30 & あまりの30° \\
\end{bktable}
\bktabcaption{くり返すものは、あまりを見ればよい。}
\end{bkcenter}

\bkivihead{反対向きに回る：マイナスの角}
時計の針は右回りに回ります。数学では、\textbf{左回り（時計の針と反対の向き）を正の向き}と決めます（4-2の極座標でも、角は左回りにはかりました）。右回りに回った角には、マイナスをつけます。2-2で「マイナスは反対の向き」と決めたのと同じ考え方です。

たとえば、右回りに120°回ることを $-120°$ と書きます。図の右のように、$-120°$ 回った所は、左回りに240°回った所と同じです。1周たして
\bkiviline{$-120°+360°=240°$}
となるからです。マイナスの角は、360°をたしていくと、0°から360°までの角に直せます。

\bkivihead{はじめの線と、回る線}
角をはかるときは、原点Oから右へのびる線（$x$ 軸の正の部分）を、はじめの線にします。これを\textbf{始線}（しせん）といいます。始線から回っていく半直線OPを、\textbf{動径}（どうけい）といいます。「径」は半径の径で、動径は「動く半径」という意味です。

30°、390°、750°、$-330°$ の動径は、どれも同じ位置です。どれも30°に、360°を何回かたしたり、ひいたりした角だからです。
\begin{bkcenter}
\begin{bktable}{colspec={Q[l,m,wd=34mm] *{4}{Q[c,m,wd=20mm]}}, rowsep=1.5pt}
360°をたす回数 $n$ & $-1$ & 0 & 1 & 2 \\
$30°+360°\times n$ & $-330°$ & 30° & 390° & 750° \\
\end{bktable}
\end{bkcenter}
このように、360°より大きい角や、マイナスの角までふくめて考えた角を、\textbf{一般角}（いっぱんかく）といいます。\bkidx{いっぱんかく}{一般角}ラジアンでは1周が $2\pi$ なので（4-2）、30°と同じ位置の角は $\dfrac{\pi}{6}+2\pi\times n$ と書けます。

\begin{seikaku}
動径が同じ位置になる角は、ある角 $\theta$ に、360°の整数倍をたした $\theta+360°\times n$（$n$ は $\cdots$、$-2$、$-1$、0、1、2、$\cdots$ の整数）です。ラジアンでは $\theta+2\pi n$ と書きます。
\end{seikaku}

\bksubsection{2}{単位円で $\sin$・$\cos$ を決めなおす}

\leavevmode\label{bk:4-8:tanien}4-6では、$\sin$ と $\cos$ を直角三角形の辺の比で決めました。直角三角形のとがった角は90°より小さいので、このままでは $\sin210°$ や $\cos750°$ は決められません。4-6と4-7では、半円の上の点の座標を使って、180°まで広げました。ここでは、それを円1周全部と、何周もする角にまで広げます。

\bkivihead{半径1の円：単位円}
原点Oを中心とする半径1の円を、\textbf{単位円}（たんいえん）といいます。\bkidx{たんいえん}{単位円}「単位」は、ここでは「1」という意味です。

単位円の上で、動径が角 $\theta$ の所にある点をPとします。Pの座標を $(x,\ y)$ とするとき、
\bkiviframe[100mm]{\sffamily\bfseries $\cos\theta$＝Pの $x$ 座標（横の位置）\qquad $\sin\theta$＝Pの $y$ 座標（高さ）}
と決めます。つまり、\textbf{単位円の上の点Pの座標は $(\cos\theta,\ \sin\theta)$} です。

\begin{bkfigure}
\begin{tikzpicture}[x=1mm, y=1mm]
  \foreach \ang/\xs in {30/0, 150/45, 210/90} {
  \begin{scope}[xshift=\xs mm]
    \draw[bkGray, line width=0.6pt, -{Stealth[length=3.5pt]}] (-20,0) -- (21,0);
    \draw[bkGray, line width=0.6pt, -{Stealth[length=3.5pt]}] (0,-19) -- (0,20);
    \draw[bkInk!45, line width=0.8pt] (0,0) circle (15);
    \draw[bkBlue, line width=1.8pt] (0,0) -- ({15*cos(\ang)},0);
    \draw[bkOrange, line width=1.8pt] ({15*cos(\ang)},0) -- (\ang:15);
    \draw[bkInk, line width=0.9pt] (0,0) -- (\ang:15);
    \draw[bkInk, line width=0.5pt] (4,0) arc (0:\ang:4);
    \fill[bkInk] (\ang:15) circle (1);
  \end{scope}}
  % 名前と座標（それぞれの図の外側に置く）
  \node[font=\sffamily\footnotesize, anchor=south west, inner sep=1pt] at ($(30:15)+(0.6,0.4)$) {P};
  \node[font=\sffamily\footnotesize] at (15:9.5) {$\theta$};
  \node[font=\sffamily\footnotesize, anchor=south east, inner sep=1pt] at ($(45,0)+(150:15)+(-0.6,0.4)$) {P};
  \node[font=\sffamily\footnotesize] at ($(45,0)+(70:6.5)$) {$\theta$};
  \node[font=\sffamily\footnotesize, anchor=north east, inner sep=1pt] at ($(90,0)+(210:15)+(-0.6,-0.4)$) {P};
  \node[font=\sffamily\footnotesize] at ($(90,0)+(70:6.5)$) {$\theta$};
  \node[font=\sffamily\footnotesize, align=center, anchor=north] at (0,-21) {$\theta=30°$\\P $(0.87,\ 0.5)$};
  \node[font=\sffamily\footnotesize, align=center, anchor=north] at (45,-21) {$\theta=150°$\\P $(-0.87,\ 0.5)$};
  \node[font=\sffamily\footnotesize, align=center, anchor=north] at (90,-21) {$\theta=210°$\\P $(-0.87,\ -0.5)$};
\end{tikzpicture}
\bkfigcaption{半径1の円（単位円）の上の点P。青い線の長さと向きが $\cos\theta$（横の位置）、オレンジの線が $\sin\theta$（高さ）。座標は小数第3位を四捨五入した。}
\end{bkfigure}

\bkivihead{三角形の決め方と同じになるわけ}
$\theta$ が0°と90°のあいだのとき（図の左）、Pから $x$ 軸に垂直な線を下ろした所をHとすると、直角三角形OHPができます。斜辺OPは単位円の半径なので1です。4-6の決め方では、
\bkiviline{$\cos\theta=\dfrac{\text{OH}}{\text{OP}}=\dfrac{\text{OH}}{1}=\text{OH}$\qquad $\sin\theta=\dfrac{\text{HP}}{\text{OP}}=\dfrac{\text{HP}}{1}=\text{HP}$}
OHはPの $x$ 座標、HPはPの $y$ 座標なので、新しい決め方と同じになります。斜辺が1だから、わり算をしても数が変わらないのです。半径を1にしたのは、このためです。

\bkivihead{90°をこえても、同じきまりで決まる}
$\theta=150°$ のとき（図のまん中）、動径OPは、$x$ 軸の左の部分と30°の角をつくっています。だから、Pから $x$ 軸に下ろした直角三角形は、30°の三角形と同じ形です。横の長さは $\cos30°\approx0.866$、高さは $\sin30°=0.5$。ただし、Pは原点より左にいるので、$x$ 座標はマイナスです。
\bkiviline{$\cos150°\approx-0.866$\qquad $\sin150°=0.5$}
これは、4-7で学んだ「鈍角の $\cos$ は負」と同じ値です。

$\theta=210°$ のとき（図の右）は、動径が $x$ 軸の左の部分から、さらに30°下に回った所です。三角形はやはり30°の形で、Pは原点より左、そして下にいます。$x$ 座標も $y$ 座標もマイナスです。
\bkiviline{$\cos210°\approx-0.866$\qquad $\sin210°=-0.5$}
{\small（0.866は、$\dfrac{\sqrt3}{2}$ を小数第3位まで表した数です。）}\par
750°なら、2周と30°なので、Pは30°の所にいます。だから、$\cos750°$ は $\cos30°$ と同じ0.866、$\sin750°$ は $\sin30°$ と同じ0.5です。

\begin{bkfigure}
\begin{tikzpicture}[x=1mm, y=1mm]
  \node[bk node, text width=25mm, font=\sffamily\footnotesize, inner xsep=1.5mm] (a) at (0,0) {① 動径をかく\\（何周かは、\\あまりで）};
  \node[bk node, text width=25mm, font=\sffamily\footnotesize, inner xsep=1.5mm] (b) at (33,0) {② $x$ 軸とのあいだのとがった角を見る};
  \node[bk node, text width=25mm, font=\sffamily\footnotesize, inner xsep=1.5mm] (c) at (66,0) {③ その角の三角形で、横とたてを出す};
  \node[bk node, text width=25mm, font=\sffamily\footnotesize, inner xsep=1.5mm] (d) at (99,0) {④ 右か左か、上か下かで符号をつける};
  \draw[bk link] (a) -- (b); \draw[bk link] (b) -- (c); \draw[bk link] (c) -- (d);
\end{tikzpicture}
\bkfigcaption{どんな角でも、$\sin$・$\cos$ はこの順で求められる。}
\end{bkfigure}

\bkivihead{三角関数という名前}
角 $\theta$ を1つ決めると、動径の位置が1つに決まり、Pの座標 $(\cos\theta,\ \sin\theta)$ も1つに決まります。$x$ を決めると $y$ が1つに決まるとき、$y$ は $x$ の関数でした（3-3）。だから、$\sin\theta$ と $\cos\theta$ は、$\theta$ の関数です。$f(\theta)=\sin\theta$ と書くこともできます。

$\sin\theta$、$\cos\theta$ と、4-6の $\tan\theta=\dfrac{\sin\theta}{\cos\theta}$ をまとめて、\textbf{三角関数}（さんかくかんすう）といいます。\bkidx{さんかくかんすう}{三角関数}$\tan\theta$ は、$\dfrac{y}{x}$、つまり動径OPの傾きです（4-6の「$\tan\theta$＝傾き」）。$\cos\theta=0$ になる90°や270°では、0でわれないので、$\tan\theta$ は決まりません。

\bkivihead{半径が $r$ の円：どの向きでも $x=r\cos\theta$、$y=r\sin\theta$}
半径 $r$ の円の上で、角 $\theta$ の所にある点Qを考えます。Qは、単位円の上の点Pを、Oを中心に $r$ 倍に拡大した点です。拡大すると、どの長さも $r$ 倍になるので（4-3）、Qの座標は、Pの座標の $r$ 倍です。
\bkiviline{Q\ $(r\cos\theta,\ r\sin\theta)$}
Qは、原点からの距離が $r$ で、向きの角が $\theta$ の点です。つまり、Qの極座標は $(r,\ \theta)$ です（4-2）。4-6では、$\theta$ が\ang{90}より小さいとき（⑧では\ang{180}まで）、極座標を $(x,\ y)$ に直す式 $x=r\cos\theta$、$y=r\sin\theta$ を作りました。$\sin$ と $\cos$ を単位円で決めなおしたので、この式が、\textbf{どの向きの点でも}そのまま使えるようになりました。
\bkiviframe[110mm]{\sffamily\bfseries 極座標 $(r,\ \theta)$ の点は\qquad $x=r\cos\theta$、\ $y=r\sin\theta$\\[2pt]{\mdseries\small（$\theta$ は何度でもよい。符号は $\cos\theta$・$\sin\theta$ が自動でつけてくれる）}}
極座標 $(2,\ 210°)$ の点なら、$x=2\cos210°=2\times(-0.866)\approx-1.73$、$y=2\sin210°=2\times(-0.5)=-1$ で、原点の左下にあります。左下にあることを、$\cos$ と $\sin$ のマイナスが教えてくれるのです。

たとえば、半径\qty{50}{m}の観覧車で、中心から見て30°の所にあるゴンドラは、中心から横に約\qty{43.3}{m}、たてに\qty{25}{m}の所にいます。
\bkiviline{横：$50\cos30°\approx50\times0.866=43.3$\qquad たて：$50\sin30°=50\times0.5=25$}
やってみようの半径\qty{10}{cm}の円でも、30°の点の高さは $10\sin30°=\qty{5}{cm}$ でした。

\bksubsection{3}{符号と、くり返し（周期）}

\bkivihead{区画ごとの符号}\label{bk:4-8:shuki}
$x$ 軸と $y$ 軸で、平面は4つの区画に分かれます。右上から左回りに、\textbf{第1象限}（しょうげん）、第2象限、第3象限、第4象限といいます。「象限」は、平面を4つに分けた1つ1つの区画のことです（軸の上の点は、どの象限にも入れません）。

\noindent\begin{minipage}[c]{0.46\linewidth}
\centering
\begin{tikzpicture}[x=0.82mm, y=0.82mm]
  \fill[bkBlue!7] (0,0) rectangle (25,25);
  \fill[bkOrange!8] (0,0) rectangle (-25,25);
  \fill[bkGreen!8] (0,0) rectangle (-25,-25);
  \fill[bkPurple!7] (0,0) rectangle (25,-25);
  \bkiviaxes{28}{28}
  \foreach \px/\py/\nm/\cs/\ss in {12.5/12.5/第1象限/＋/＋, -12.5/12.5/第2象限/－/＋, -12.5/-12.5/第3象限/－/－, 12.5/-12.5/第4象限/＋/－} {
    \node[font=\sffamily\footnotesize, align=center, text=bkInk] at (\px,\py) {\textbf{\nm}\\$\cos$ が \cs\\$\sin$ が \ss};}
\end{tikzpicture}
\end{minipage}\hfill
\begin{minipage}[c]{0.5\linewidth}
\raggedright
覚えなくても、Pの場所を見ればわかります。
\begin{itemize}[leftmargin=1.2\zw, itemsep=2pt, topsep=2pt, label={\color{\bkTT}●}]
\item $\cos\theta$ は横の位置。Pが $y$ 軸より右なら正、左なら負。
\item $\sin\theta$ は高さ。Pが $x$ 軸より上なら正、下なら負。
\end{itemize}
\end{minipage}\par\vspace{4pt}
たとえば250°は、180°と270°のあいだなので第3象限です。Pは左下にいるので、$\cos250°$ も $\sin250°$ も負です。電卓では $\cos250°=-0.342$、$\sin250°=-0.940$ です。

\bkivihead{1周するともとにもどる：周期}
$\theta$ に360°をたすと、動径は1周して同じ位置にもどります。Pも同じ点なので、座標も同じです。
\bkiviline{$\sin(\theta+360°)=\sin\theta$\qquad $\cos(\theta+360°)=\cos\theta$}
電卓で $\sin390°$ を計算すると、$\sin30°$ と同じ0.5になります。$\sin$ と $\cos$ の値は、360°ごとに同じ変わり方をくり返します。このように、同じ変わり方をくり返すとき、1回分のはばを\textbf{周期}（しゅうき）といいます。\bkidx{しゅうき}{周期}$\sin\theta$ と $\cos\theta$ の周期は360°、ラジアンでは $2\pi$ です。

\bkivihead{3つの「うつしかた」}
角を変えると、Pがどこへ動くかを、図で考えます。4-1で学んだ「裏返す（線対称）」と「回す（点対称）」が使えます。

\noindent\begin{minipage}[c]{0.46\linewidth}
\centering
\begin{tikzpicture}[x=1mm, y=1mm]
  \bkiviaxes{25}{23}
  \draw[bkInk!45, line width=0.8pt] (0,0) circle (17);
  \draw[bkGray, dashed, line width=0.5pt] (30:17) -- (150:17) -- (210:17) -- (330:17) -- cycle;
  \foreach \a in {30,150,210,330} {\draw[bkInk, line width=0.8pt] (0,0) -- (\a:17); \fill[bkInk] (\a:17) circle (1);}
  \fill[bkOrange] (30:17) circle (1.3);
  \node[font=\sffamily\footnotesize, anchor=south west, inner sep=1pt, align=left] at ($(30:17)+(0.5,0.6)$) {P\enspace $\theta$\\$(x,\ y)$};
  \node[font=\sffamily\footnotesize, anchor=south east, inner sep=1pt, align=right] at ($(150:17)+(-0.5,0.6)$) {$180°-\theta$\\$(-x,\ y)$};
  \node[font=\sffamily\footnotesize, anchor=north east, inner sep=1pt, align=right] at ($(210:17)+(-0.5,-0.6)$) {$(-x,\ -y)$\\$\theta+180°$};
  \node[font=\sffamily\footnotesize, anchor=north west, inner sep=1pt, align=left] at ($(330:17)+(0.5,-0.6)$) {$(x,\ -y)$\\$-\theta$};
\end{tikzpicture}
\end{minipage}\hfill
\begin{minipage}[c]{0.5\linewidth}
\begin{itemize}[leftmargin=1.2\zw, itemsep=3pt, topsep=0pt, label={\color{\bkTT}●}]
\item \textbf{$-\theta$}\quad 始線から反対向きに同じだけ回る。$x$ 軸で裏返した位置なので、横はそのまま、高さは反対。
\item \textbf{$\theta+180°$}\quad 半周回る。原点について点対称（4-1：点対称は180°回すこと）なので、横も高さも反対。
\item \textbf{$180°-\theta$}\quad $y$ 軸で裏返した位置（4-6・4-7でも使った）。横は反対、高さはそのまま。
\end{itemize}
\end{minipage}

\begin{bkcenter}
\begin{bktable}{colspec={Q[c,m,wd=14mm] Q[c,m,wd=24mm] Q[c,m,wd=37mm] Q[c,m,wd=37mm]}, rowsep=2pt, colsep=3pt, cells={font=\small}, row{1}={font=\sffamily\small}}
角 & うつしかた & $\cos$ & $\sin$ \\
$-\theta$ & $x$ 軸で裏返す & $\cos(-\theta)=\cos\theta$ & $\sin(-\theta)=-\sin\theta$ \\
$\theta+180°$ & 原点のまわりに半周 & $\cos(\theta+180°)=-\cos\theta$ & $\sin(\theta+180°)=-\sin\theta$ \\
$180°-\theta$ & $y$ 軸で裏返す & $\cos(180°-\theta)=-\cos\theta$ & $\sin(180°-\theta)=\sin\theta$ \\
\end{bktable}
\bktabcaption{覚えるより、図をかいて考えるほうが確か。}
\end{bkcenter}
数で確かめましょう。$\theta=30°$ なら、$\cos(-30°)$ は $\cos30°$ と同じ約0.866、$\sin(-30°)$ は $-\sin30°$ と同じ $-0.5$ です。$\theta+180°=210°$ では、②で求めたとおり、$\cos210°$ は約 $-0.866$、$\sin210°$ は $-0.5$ で、どちらも30°の値の反対です。

\bksubsection{4}{$\sin^2\theta+\cos^2\theta=1$：単位円で三平方の定理}

\leavevmode\label{bk:4-8:sanheihou}単位円の上の点Pは、どこにいても、原点Oからの距離が1です（半径だから）。Pの座標は $(\cos\theta,\ \sin\theta)$ なので、三平方の定理（4-4）で OP の長さの2乗を計算すると、
\bkiviline{$\text{OP}^2=(\cos\theta)^2+(\sin\theta)^2=1^2=1$}
4-4の円の式 $x^2+y^2=r^2$ に、$x=\cos\theta$、$y=\sin\theta$、$r=1$ を入れたものと同じです。
\par\vspace{6pt}
\noindent\begin{minipage}[c]{0.6\linewidth}
\raggedright
Pが第2・第3象限にいて、$\cos\theta$ がマイナスのときも、この式は成り立ちます。三角形の横の長さは「$\cos\theta$ のマイナスをとった数」ですが、マイナスの数も2乗すると正になるので、$(\cos\theta)^2$ は横の長さの2乗と同じになるからです。$\sin\theta$ がマイナスのときも同じです。

$(\sin\theta)^2$ は、ふつう $\sin^2\theta$ と書きます。2を $\sin$ のすぐ右上に書きます。
\end{minipage}\hfill
\begin{minipage}[c]{0.36\linewidth}
\centering
\begin{tikzpicture}[x=0.85mm, y=0.85mm]
  \bkiviaxes{23}{22}
  \draw[bkInk!45, line width=0.8pt] (0,0) circle (18);
  \draw[bkBlue, line width=1.8pt] (0,0) -- ({18*cos(210)},0);
  \draw[bkOrange, line width=1.8pt] ({18*cos(210)},0) -- (210:18);
  \draw[bkInk, line width=0.9pt] (0,0) -- (210:18);
  \fill[bkInk] (210:18) circle (1);
  \node[font=\sffamily\footnotesize, anchor=south, text=bkBlue!80!black, inner sep=1pt] at (-7.8,0.8) {横};
  \node[font=\sffamily\footnotesize, text=bkOrange!80!black, inner sep=1pt] at (-12,-2.4) {たて};
  \node[font=\sffamily\footnotesize, anchor=north west, inner sep=1pt] at (-6,-5.2) {1};
  \node[font=\sffamily\footnotesize, anchor=north east, inner sep=1pt] at ($(210:18)+(-0.5,-0.5)$) {P};
\end{tikzpicture}
\end{minipage}
\bkiviframe[60mm]{\large $\sin^2\theta+\cos^2\theta=1$}
$\sin^2\theta$ は「$\sin\theta$ を2乗した数」です。$\sin\theta^2$ と書くと、「$\theta^2$ の $\sin$」という、まったく別のものになってしまいます。

\begin{bkcenter}
\begin{bktable}{colspec={Q[c,m,wd=16mm] Q[c,m,wd=22mm] Q[c,m,wd=22mm] Q[c,m,wd=48mm]}, rowsep=1.5pt}
$\theta$ & $\cos\theta$ & $\sin\theta$ & $\cos^2\theta+\sin^2\theta$ \\
30° & 0.866 & 0.5 & $0.75+0.25=1$ \\
135° & $-0.707$ & 0.707 & $0.5+0.5=1$ \\
210° & $-0.866$ & $-0.5$ & $0.75+0.25=1$ \\
\end{bktable}
\bktabcaption{どの象限の角でも、2乗してたすと1。（$0.866^2=0.75$、$0.707^2=0.5$ は、$\bigl(\frac{\sqrt3}{2}\bigr)^2=\frac34$、$\bigl(\frac{\sqrt2}{2}\bigr)^2=\frac24$ を小数にしたもの）}
\end{bkcenter}

\bkivihead{片方から、もう片方を求める}
この式を使うと、$\sin\theta$ がわかれば $\cos\theta$ が、$\cos\theta$ がわかれば $\sin\theta$ が求められます。たとえば、$\sin\theta=0.6$ で、$\theta$ が第2象限の角のとき、
\bkiviline{$\cos^2\theta=1-\sin^2\theta=1-0.6^2=1-0.36=0.64$}
2乗して0.64になる数は、0.8と $-0.8$ の2つです。単位円で見ると、高さ0.6の点は、単位円の右側と左側に1つずつあります。$\theta$ は第2象限（左上）の角なので、$\cos\theta$ は負で、$\cos\theta=-0.8$ です。

\bksubsection{5}{90°回すと $(x,\ y)\to(-y,\ x)$}

\leavevmode\label{bk:4-8:kaiten}ロボットのうでや、時計の針を、原点のまわりに左回りに90°回すと、先の点はどこへ動くでしょう。点P $(3,\ 1)$ を、原点Oのまわりに左回りに90°回した点P$'$（ピー・ダッシュ）を、2通りで求めます。

\bkivihead{見方1：方眼の上で、道すじごと回す}
\noindent\begin{minipage}[c]{0.42\linewidth}
\centering
\begin{tikzpicture}[x=6.3mm, y=6.3mm]
  \draw[bkLine, line width=0.4pt] (-2.4,-0.6) grid (3.6,3.6);
  \draw[bkGray, line width=0.6pt, -{Stealth[length=4pt]}] (-2.6,0) -- (3.8,0) node[anchor=west, font=\sffamily\footnotesize, text=bkInk] {$x$};
  \draw[bkGray, line width=0.6pt, -{Stealth[length=4pt]}] (0,-0.8) -- (0,3.8) node[anchor=south, font=\sffamily\footnotesize, text=bkInk] {$y$};
  \draw[bkBlue, line width=1.8pt, -{Stealth[length=5pt]}] (0,0) -- (2.92,0);
  \draw[bkOrange, line width=1.8pt, -{Stealth[length=5pt]}] (3,0) -- (3,0.92);
  \draw[bkBlue, line width=1.8pt, -{Stealth[length=5pt]}] (0,0) -- (0,2.92);
  \draw[bkOrange, line width=1.8pt, -{Stealth[length=5pt]}] (0,3) -- (-0.92,3);
  \draw[bkGray, line width=0.7pt, dash pattern=on 2pt off 1.5pt, -{Stealth[length=4pt]}] ({sqrt(10)*cos(24)},{sqrt(10)*sin(24)}) arc (24:{90+14}:{sqrt(10)});
  \fill[bkInk] (3,1) circle (0.12) (-1,3) circle (0.12);
  \node[font=\sffamily\footnotesize, anchor=west] at (3.12,1.05) {P $(3,\ 1)$};
  \node[font=\sffamily\footnotesize, anchor=east] at (-1.12,3.0) {P$'$};
  \node[font=\sffamily\footnotesize, text=bkBlue!80!black, anchor=north, fill=white, inner sep=1pt] at (1.5,-0.16) {右へ3};
  \node[font=\sffamily\footnotesize, text=bkBlue!80!black, anchor=west, fill=white, inner sep=1pt] at (0.18,1.5) {上へ3};
  \node[font=\sffamily\footnotesize, text=bkOrange!80!black, anchor=south, fill=white, inner sep=1pt] at (-0.5,3.18) {左へ1};
  \node[font=\sffamily\footnotesize, text=bkOrange!80!black, anchor=east, fill=white, inner sep=1pt] at (2.82,0.5) {上へ1};
  \node[font=\sffamily\footnotesize, text=bkGray, fill=white, inner sep=1pt] at (62:3.8) {90°};
\end{tikzpicture}
\end{minipage}\hfill
\begin{minipage}[c]{0.55\linewidth}
\raggedright
Pへは、Oから「右へ3（青）、上へ1（オレンジ）」と進めば着きます。この道すじを、まるごと左回りに90°回します。
\begin{itemize}[leftmargin=1.2\zw, itemsep=2pt, topsep=2pt, label={\color{\bkTT}●}]
\item 「右へ3」は、90°回ると「上へ3」
\item 「上へ1」は、90°回ると「左へ1」
\end{itemize}
だから、P$'$ は「左へ1、上へ3」の所で、P$'$ $(-1,\ 3)$ です。
\end{minipage}\par\vspace{4pt}
長さを確かめます。OP$^2=3^2+1^2=10$、OP$'^2=(-1)^2+3^2=10$ で、原点からの距離は変わっていません（4-4）。回しても、長さは変わらないからです。

\bkivihead{見方2：直角三角形をまるごと90°回す}
こんどは、点P $(x,\ y)$ を一般に考えます（下の図は、$x$ も $y$ も正のときの図です）。Pから $x$ 軸に垂直な線を下ろすと、横が $x$（青）、たてが $y$（オレンジ）の直角三角形ができます。この三角形を、Oを中心に、\textbf{まるごと90°回します}。

\begin{bkfigure}
\begin{tikzpicture}[x=1mm, y=1mm]
  \def\R{30}
  \draw[bkGray, line width=0.6pt, -{Stealth[length=4pt]}] (-22,0) -- (38,0) node[anchor=west, font=\sffamily\footnotesize] {$x$};
  \draw[bkGray, line width=0.6pt, -{Stealth[length=4pt]}] (0,-4) -- (0,36) node[anchor=south, font=\sffamily\footnotesize] {$y$};
  \draw[bkInk!35, line width=0.8pt] (-5:\R) arc (-5:130:\R);
  % もとの三角形（横 x は青、たて y はオレンジ）
  \fill[bkBlue!8] (0,0) -- (25.98,0) -- (25.98,15) -- cycle;
  \draw[bkBlue, line width=1.8pt] (0,0) -- (25.98,0);
  \draw[bkOrange, line width=1.8pt] (25.98,0) -- (25.98,15);
  \draw[bkInk, line width=0.9pt] (0,0) -- (25.98,15);
  \fill[bkInk] (25.98,15) circle (1);
  \node[font=\sffamily\footnotesize, anchor=north, text=bkBlue!80!black] at (13,-1.2) {$x$};
  \node[font=\sffamily\footnotesize, anchor=west, text=bkOrange!80!black] at (26.8,7.5) {$y$};
  \node[font=\sffamily\footnotesize, anchor=south west, inner sep=1pt] at (26.6,15.4) {P $(x,\ y)$};
  % 90°回した三角形（横 y はオレンジ、たて x は青）
  \fill[bkOrange!8] (0,0) -- (0,25.98) -- (-15,25.98) -- cycle;
  \draw[bkBlue, line width=1.8pt] (0,0) -- (0,25.98);
  \draw[bkOrange, line width=1.8pt] (0,25.98) -- (-15,25.98);
  \draw[bkInk, line width=0.9pt] (0,0) -- (-15,25.98);
  \fill[bkInk] (-15,25.98) circle (1);
  \node[font=\sffamily\footnotesize, anchor=west, text=bkBlue!80!black] at (0.8,13) {$x$};
  \node[font=\sffamily\footnotesize, anchor=north, text=bkOrange!80!black] at (-5.5,25.2) {$y$};
  \node[font=\sffamily\footnotesize, anchor=east, inner sep=1pt] at (-16,25.98) {P$'$};
  \draw[bkGray, line width=0.7pt, dash pattern=on 2pt off 1.5pt, -{Stealth[length=4pt]}] (34:\R) arc (34:116:\R);
  \node[font=\sffamily\footnotesize, text=bkGray, fill=white, inner sep=1pt] at (75:{\R+3.5}) {90°};
  % 右の説明
  \node[font=\sffamily\small, anchor=north west, align=left, text width=62mm] at (46,34) {%
  \textcolor{bkBlue!80!black}{青い辺}の長さは、どちらも $x$。\\
  \textcolor{bkOrange!80!black}{オレンジの辺}の長さは、どちらも $y$。\\[4pt]
  回した三角形では、\\
  \quad 青い辺は、\textbf{たて}にO から上へ $x$\\
  \quad オレンジの辺は、\textbf{横}に左へ $y$\\[4pt]
  だから、P$'$ の座標は $(-y,\ x)$};
\end{tikzpicture}
\bkfigcaption{三角形（青い面）を90°回すと、オレンジの面の三角形になる。辺の長さはそのまま、たてと横が入れかわり、横は左向きになる。}
\end{bkfigure}

\bkiviframe[100mm]{左回りに90°回すと\qquad $(x,\ y)\ \to\ (-y,\ x)$}
見方1の P $(3,\ 1)$ で確かめると、$(-1,\ 3)$ になり、合っています。この90°回す式は、Pがどの象限にいても成り立ちます（$x$ や $y$ がマイナスのときは、向きも反対になって、うまく合います）。

\bkivihead{$\sin$・$\cos$ で書くと}
単位円の上の点P $(\cos\theta,\ \sin\theta)$ を、90°回します。回した点P$'$ は、動径が $\theta+90°$ の所にあるので、$x$ 座標は $\cos(\theta+90°)$、$y$ 座標は $\sin(\theta+90°)$ です。いっぽう、90°回す式 $(x,\ y)\to(-y,\ x)$ からは、$x$ 座標が $-\sin\theta$、$y$ 座標が $\cos\theta$ です。2つは同じ点なので、
\bkiviframe[100mm]{$\cos(\theta+90°)=-\sin\theta$\qquad $\sin(\theta+90°)=\cos\theta$}
$\theta=30°$ で確かめます。$\theta+90°=120°$ です。
\begin{align*}
\cos120°&=-0.5 & -\sin30°&=-0.5\\
\sin120°&=0.866 & \cos30°&=0.866
\end{align*}
左右が等しく、合っています。

90°回すことを2回くり返すと、$(x,\ y)\to(-y,\ x)\to(-x,\ -y)$ で、③の「180°回すと、横も高さも反対」と同じになります。この90°回す見方は、5-15で回る点の速度の向きを調べるときと、4-15の極座標で使います。

\bksubsection{6}{グラフ：回る点の影は、波になる}

\leavevmode\label{bk:4-8:graph}単位円の上を、Pが左回りに回ります。Pの高さ $\sin\theta$ は、真横から光を当てたときの、かべの上の影の高さです。横軸に角 $\theta$、たて軸に影の高さをとって、グラフにかきます。

\begin{bkfigure}
\begin{tikzpicture}[x=1mm, y=1mm]
  % 左：単位円（半径 16 mm）、右：y=sinθ のグラフ（360° を 80 mm、1 を 16 mm）
  \def\R{16}
  \draw[bkGray, line width=0.6pt] (-19,0) -- (19,0);
  \draw[bkGray, line width=0.6pt] (0,-19) -- (0,19);
  \draw[bkInk!45, line width=0.8pt] (0,0) circle (\R);
  \foreach \a in {0,30,...,330} {\fill[bkBlue] (\a:\R) circle (0.8);}
  \draw[bkGray, line width=0.6pt, -{Stealth[length=4pt]}] (25,0) -- (109,0) node[anchor=west, font=\sffamily\footnotesize] {$\theta$};
  \draw[bkGray, line width=0.6pt, -{Stealth[length=4pt]}] (25,-20) -- (25,21) node[anchor=south, font=\sffamily\footnotesize] {$y$};
  \foreach \k in {1,2,3,4} {\draw[bkGray, line width=0.5pt] ({25+20*\k},-0.8) -- ({25+20*\k},0.8);}
  \node[font=\sffamily\footnotesize, anchor=south west, inner sep=1pt] at (65.4,1.0) {180°};
  \node[font=\sffamily\footnotesize, anchor=south east, inner sep=1pt] at (104.6,1.0) {360°};
  \foreach \k/\lab in {1/90°} {\node[font=\sffamily\footnotesize, anchor=north, inner sep=1pt] at ({25+20*\k},-1.2) {\lab};}
  \foreach \k/\lab in {3/270°} {\node[font=\sffamily\footnotesize, anchor=south, inner sep=1pt] at ({25+20*\k},1.2) {\lab};}
  \draw[bkGray, line width=0.5pt] (24.2,16) -- (25.8,16) (24.2,-16) -- (25.8,-16);
  \node[font=\sffamily\footnotesize, anchor=east] at (23.8,16) {1};
  \node[font=\sffamily\footnotesize, anchor=east] at (23.8,-16) {$-1$};
  \draw[bkBlue, line width=1.6pt, domain=0:360, samples=120, smooth, variable=\t] plot ({25+\t/4.5},{\R*sin(\t)});
  \foreach \a in {0,30,...,360} {\fill[bkBlue] ({25+\a/4.5},{\R*sin(\a)}) circle (0.8);}
  % 影の線（30°・60°・90°）。左半分の点から引くと円や曲線と重なるので、右上の3つだけ
  \foreach \a in {30,60,90} {\draw[bkOrange, line width=0.6pt, dash pattern=on 2pt off 1.5pt] (\a:\R) -- ({25+\a/4.5},{\R*sin(\a)});}
  \fill[bkOrange] (30:\R) circle (1.2) (60:\R) circle (1.2) (90:\R) circle (1.2);
\end{tikzpicture}
\bkfigcaption{単位円の上の点（左）の高さを、角ごとに右へならべると、$y=\sin\theta$ のグラフになる。オレンジの点線は、30°・60°・90°の点の高さ。}
\end{bkfigure}

やってみようでつないだ形と同じ、山と谷のある波になりました。このグラフからわかることを書き出します。
\begin{itemize}[leftmargin=1.2\zw, itemsep=2pt, topsep=2pt, label={\color{\bkTT}●}]
\item いちばん高いのは90°で1、いちばん低いのは270°で $-1$。$\sin\theta$ の値は、$-1$ から1までのあいだにある（Pは半径1の円の上にいるから）。
\item 0°、180°、360°で0（Pが $x$ 軸の上にいるから）。
\item 360°から先は、同じ形をくり返す（周期360°）。マイナスの角でも、左へ同じ形が続く。
\end{itemize}

\bkivihead{影が、上と下ではゆっくり動くわけ}
グラフは、0°や180°のあたりでは急で、90°や270°のあたりでは平らです。同じ30°回るあいだの高さの変わり方を、数でくらべます。
\begin{bkcenter}
\begin{bktable}{colspec={Q[c,m,wd=32mm] Q[c,m,wd=44mm] Q[c,m,wd=28mm]}, rowsep=1.5pt}
回るあいだ & 高さ & 高さの変化 \\
0°→30° & 0 → 0.5 & 0.5 \\
30°→60° & 0.5 → 0.866 & 0.366 \\
60°→90° & 0.866 → 1 & 0.134 \\
\end{bktable}
\end{bkcenter}
同じ30°ずつ回っても、高さの変化は0.5、0.366、0.134と小さくなっていきます。Pが真横（0°）を通るときは、Pはまっすぐ上に動いているので、高さがどんどん変わります。Pがいちばん上（90°）の近くでは、Pはほとんど横に動いているので、高さはあまり変わりません。これが、\bkhatenaref の「影が上と下ではゆっくり、まん中で速く動く」わけです。

\bkivihead{$y=\cos\theta$ のグラフ}
Pの横の位置 $\cos\theta$ は、真上から光を当てたときの、床の上の影の位置です。$y=\cos\theta$ のグラフは、0°で1から始まり、90°で0、180°で $-1$ になります。
\begin{bkfigure}
\begin{tikzpicture}
\begin{axis}[bkplot, width=124mm, height=46mm, xmin=-180, xmax=600, ymin=-1.3, ymax=1.3, xtick={-90,90,180,270,360,450}, xticklabels={}, ytick={-1,1}, yticklabels={}, xlabel={$\theta$}, ylabel={$y$}]
  \foreach \v/\l in {-90/{$-90°$},90/90°,180/180°,270/270°,360/360°,450/450°} {\edef\tmp{\noexpand\node[font=\noexpand\sffamily\noexpand\footnotesize, anchor=north, inner sep=1pt] at (axis cs:\v,-1.32) {\unexpanded\expandafter{\l}};}\tmp}
  \node[font=\sffamily\footnotesize, anchor=east, inner sep=1pt] at (axis cs:-182,1) {1};
  \node[font=\sffamily\footnotesize, anchor=east, inner sep=1pt] at (axis cs:-182,-1) {$-1$};
  \addplot[bkcurve, domain=-180:540, samples=181] {sin(x)};
  \addplot[bkcurve2, domain=-180:540, samples=181] {cos(x)};
  \draw[bkGray, line width=0.8pt, -{Stealth[length=4pt]}] (axis cs:90,1.12) -- (axis cs:0,1.12);
  \node[font=\sffamily\footnotesize, anchor=south, inner sep=1pt] at (axis cs:45,1.14) {90°};
  \node[font=\sffamily\footnotesize, text=bkBlue, anchor=west, inner sep=1pt] at (axis cs:548,0.4) {$\sin\theta$};
  \node[font=\sffamily\footnotesize, text=bkOrange!80!black, anchor=west, inner sep=1pt] at (axis cs:548,-0.8) {$\cos\theta$};
\end{axis}
\end{tikzpicture}
\bkfigcaption{$y=\sin\theta$（青）と $y=\cos\theta$（オレンジ）。$\cos\theta$ の山（0°）は、$\sin\theta$ の山（90°）より90°早く来る。}
\end{bkfigure}
$\cos\theta$ のグラフは、$\sin\theta$ のグラフと同じ形を、左へ90°ずらしたものです。⑤で作った $\sin(\theta+90°)=\cos\theta$ が、このことを表しています。$\theta=0°$ のとき、$\cos\theta$ は、$\sin$ のグラフの $\theta=90°$ の値（山の1）と同じです。
\begin{bkcenter}
\begin{bktable}{colspec={Q[l,m,wd=28mm] *{4}{Q[c,m,wd=17mm]}}, rowsep=1.5pt}
$\theta$ & 0° & 30° & 60° & 90° \\
$\cos\theta$ & 1 & 0.866 & 0.5 & 0 \\
$\theta+90°$ & 90° & 120° & 150° & 180° \\
$\sin(\theta+90°)$ & 1 & 0.866 & 0.5 & 0 \\
\end{bktable}
\end{bkcenter}
$\sin\theta$ のグラフの形を、\textbf{正弦曲線}（せいげんきょくせん）、英語でサインカーブといいます。$\sin$ を日本語で正弦（せいげん）、$\cos$ を余弦（よげん）といいます。

\bksubsection{7}{波をのばす・ずらす}

\leavevmode\label{bk:4-8:idou}3-9では、放物線 $y=a(x-p)^2+q$ が、$y=ax^2$ を右へ $p$、上へ $q$ 動かしたものだと学びました。三角関数のグラフでも、同じことができます。4つの変え方を、1つずつ見ます。どの図も、灰色の線がもとの $y=\sin\theta$ です。

\begin{bkfigure}
\begin{tikzpicture}
\pgfplotsset{bkivisml/.style={bkplot, width=61mm, height=38mm, xmin=0, xmax=400, ymin=-2.3, ymax=2.5, xtick={90,180,270,360}, xticklabels={}, ytick={-2,-1,1,2}, tick label style={font=\sffamily\footnotesize}, title style={font=\sffamily\small, yshift=-1mm},
  extra description/.code={\foreach \v/\l in {90/90°,180/180°,270/270°,360/360°} {\node[font=\sffamily\footnotesize, anchor=north, inner sep=1pt] at (axis cs:\v,-2.32) {\l};}}}}
\begin{axis}[bkivisml, name=A, title={(a)\enspace $y=2\sin\theta$（たてに2倍）}]
  \addplot[bkthin, domain=0:380, samples=96] {sin(x)};
  \addplot[bkcurve, domain=0:380, samples=96] {2*sin(x)};
\end{axis}
\begin{axis}[bkivisml, at={($(A.east)+(8mm,0)$)}, anchor=west, title={(b)\enspace $y=\sin\theta+1$（上へ1）}]
  \addplot[bkthin, domain=0:380, samples=96] {sin(x)};
  \addplot[bkcurve, domain=0:380, samples=96] {sin(x)+1};
\end{axis}
\begin{axis}[bkivisml, at={($(A.south)+(0,-14mm)$)}, anchor=north, name=C, title={(c)\enspace $y=\sin(\theta-60°)$（右へ60°）}]
  \addplot[bkthin, domain=0:380, samples=96] {sin(x)};
  \addplot[bkcurve, domain=0:380, samples=96] {sin(x-60)};
  \draw[bkOrange!85!black, line width=0.9pt, -{Stealth[length=4pt]}] (axis cs:90,1.3) -- (axis cs:150,1.3);
  \node[font=\sffamily\footnotesize, anchor=south, inner sep=1pt] at (axis cs:120,1.35) {60°};
\end{axis}
\begin{axis}[bkivisml, at={($(C.east)+(8mm,0)$)}, anchor=west, title={(d)\enspace $y=\sin2\theta$（周期が半分）}]
  \addplot[bkthin, domain=0:380, samples=96] {sin(x)};
  \addplot[bkcurve, domain=0:380, samples=190] {sin(2*x)};
\end{axis}
\end{tikzpicture}
\bkfigcaption{青い線が新しいグラフ、灰色の線がもとの $y=\sin\theta$。}
\end{bkfigure}

\begin{itemize}[leftmargin=1.2\zw, itemsep=3pt, topsep=2pt, label={\color{\bkTT}●}]
\item \textbf{(a)\enspace$y=2\sin\theta$}\\
どの値も2倍なので、山は2、谷は $-2$。半径2の円を回る点の高さ $2\sin\theta$ と同じです（②の $r\sin\theta$）。中心の線から山までの高さを、\textbf{振幅}（しんぷく）といいます。$y=2\sin\theta$ の振幅は2です。
\item \textbf{(b)\enspace$y=\sin\theta+1$}\\
どの値にも1をたすので、グラフ全体が上へ1動きます。3-9の $+q$ と同じです。
\item \textbf{(c)\enspace$y=\sin(\theta-60°)$}\\
$y=\sin\theta$ が $\theta=0°$ でしていたことを、$y=\sin(\theta-60°)$ は、かっこの中が0になる $\theta=60°$ でします。どのできごとも60°おそく起こるので、グラフは\textbf{右へ}60°動きます。3-9で $y=(x-3)^2$ の頂点が $x=3$（右）に動いたのと同じです。
\item \textbf{(d)\enspace$y=\sin2\theta$}\\
$\theta$ が180°になると、$2\theta$ が360°になって、1周分の波が終わります。2倍の速さで回る点の高さと同じで、周期は $360°\div2=180°$ です。
\end{itemize}

(c)を数で確かめます。$y=\sin\theta$ は $\theta=0°$、30°、90°で0、0.5、1でした。$y=\sin(\theta-60°)$ に $\theta=60°$、90°、150°を入れると、$\sin0°=0$、$\sin30°=0.5$、$\sin90°=1$。同じ値が、60°ずつおくれて出てきます。

\begin{bkcenter}
\begin{bktable}{colspec={Q[c,m,wd=30mm] Q[l,m,wd=46mm] Q[c,m,wd=30mm]}, rowsep=2pt, row{1}={font=\sffamily\small}}
式 & グラフの変わり方 & 3-9の放物線では \\
$y=a\sin\theta$ & たてに $a$ 倍（振幅 $a$） & $y=ax^2$ \\
$y=\sin\theta+q$ & 上へ $q$ & $y=x^2+q$ \\
$y=\sin(\theta-p)$ & 右へ $p$ & $y=(x-p)^2$ \\
$y=\sin k\theta$ & 周期が $360°\div k$ & ― \\
\end{bktable}
\bktabcaption{三角関数のグラフの動かし方。$\cos$ でも同じ。}
\end{bkcenter}

\bksubsection{8}{$\sin\theta=0.5$ になる角は、1周に2つ}

\leavevmode\label{bk:4-8:futatsu}4-6では、電卓の $\sin^{-1}$ で「$\sin\theta=0.5$ になる角」を求めると、30°が出てきました。でも、一般角まで考えると、$\sin\theta=0.5$ になる角はほかにもあります。2通りの見方で探します。

\begin{bkfigure}
\begin{tikzpicture}[x=1mm, y=1mm]
  \def\R{14}
  \bkiviaxes{19}{19}
  \draw[bkInk!45, line width=0.8pt] (0,0) circle (\R);
  \draw[bkOrange, line width=1pt] (-18,7) -- (18,7);
  \foreach \a in {30,150} {\draw[bkInk, line width=0.8pt] (0,0) -- (\a:\R); \fill[bkInk] (\a:\R) circle (1.1);}
  \node[font=\sffamily\footnotesize, anchor=south west, inner sep=1pt] at ($(30:\R)+(0.4,0.6)$) {30°};
  \node[font=\sffamily\footnotesize, anchor=south east, inner sep=1pt] at ($(150:\R)+(-0.4,0.6)$) {150°};
  \begin{scope}[xshift=30mm]
  \draw[bkGray, line width=0.6pt, -{Stealth[length=4pt]}] (0,0) -- (80,0) node[anchor=west, font=\sffamily\footnotesize] {$\theta$};
  \draw[bkGray, line width=0.6pt, -{Stealth[length=4pt]}] (0,-17) -- (0,18) node[anchor=south, font=\sffamily\footnotesize] {$y$};
  \draw[bkBlue, line width=1.6pt, domain=0:720, samples=200, smooth, variable=\t] plot ({\t/9.6},{\R*sin(\t)});
  \draw[bkOrange, line width=1pt] (0,7) -- (77,7);
  \foreach \a in {30,150,390,510} {\fill[bkInk] ({\a/9.6},7) circle (1); \draw[bkGray, line width=0.5pt, dash pattern=on 1.5pt off 1.5pt] ({\a/9.6},7) -- ({\a/9.6},0);}
  \foreach \a/\lab in {30/30°,150/150°,390/390°,510/510°} {\node[font=\sffamily\footnotesize, anchor=north, fill=white, inner sep=1pt] at ({\a/9.6},-1.2) {\lab};}
  \draw[bkGray, line width=0.5pt] (-0.8,14) -- (0.8,14) (-0.8,-14) -- (0.8,-14);
  \node[font=\sffamily\footnotesize, anchor=east] at (-1.2,14) {1};
  \node[font=\sffamily\footnotesize, anchor=east, text=bkOrange!80!black] at (-1.2,7) {0.5};
  \node[font=\sffamily\footnotesize, anchor=east] at (-1.2,-14) {$-1$};
  \end{scope}
\end{tikzpicture}
\bkfigcaption{左：高さ0.5の横線（オレンジ）は、単位円と2点で交わる。右：$y=\sin\theta$ のグラフと横線 $y=0.5$ は、1周（360°）ごとに2つずつ交わる。}
\end{bkfigure}

\begin{itemize}[leftmargin=1.2\zw, itemsep=3pt, topsep=2pt, label={\color{\bkTT}●}]
\item \textbf{単位円で}\quad $\sin\theta$ は高さなので、高さ0.5の横線と単位円が交わる点を探します。交わる点は、右に1つ（30°）、左に1つ（150°）あります。左の点は、右の点を $y$ 軸で裏返した位置で、$180°-30°=150°$ です（③の表）。
\item \textbf{グラフで}\quad $y=\sin\theta$ のグラフと、横線 $y=0.5$ が交わる所を探します。0°から360°のあいだに2つ（30°と150°）。さらに、360°ずつ先に、390°、510°、……と続きます。
\end{itemize}
電卓の $\sin^{-1}$ は、1つ（30°）しか教えてくれません。もう1つは、自分で $180°-30°$ と考えます。

\bkivihead{$\cos\theta$ のときは}
$\cos\theta$ は横の位置なので、$\cos\theta=0.5$ になる角は、たての線 $x=0.5$ と単位円が交わる所です。上に1つ（60°）、下に1つあります。下の点は、上の点を $x$ 軸で裏返した位置で、$-60°$、つまり $360°-60°=300°$ です（③の $\cos(-\theta)=\cos\theta$）。
\begin{bkcenter}
\begin{bktable}{colspec={Q[c,m,wd=26mm] Q[c,m,wd=30mm] Q[c,m,wd=44mm]}, rowsep=2pt, row{1}={font=\sffamily\small}}
式 & 電卓で出る角 $\theta_1$ & もう1つの角（0°〜360°） \\
$\sin\theta=k$ & $\sin^{-1}k$ & $180°-\theta_1$ \\
$\cos\theta=k$ & $\cos^{-1}k$ & $360°-\theta_1$ \\
\end{bktable}
\bktabcaption{$\theta_1$（シータ・いち）は、電卓で出た角。マイナスの角が出たときは、360°をたして直す。}
\end{bkcenter}
$\sin\theta$ がマイナスになる角も、同じように考えます。$\sin\theta=-0.5$ なら、高さ $-0.5$ の横線と交わる2点で、210°と330°です（30°の三角形が、左下と右下にできる）。

\bksubsection{9}{回る点を時間で表す：角速度 $\omega$}

\leavevmode\label{bk:4-8:omega}回転台が、4秒で1周する速さで、同じように回っています。回転台のふちの点は、$t$ 秒でどれだけの角を回るでしょう。
\begin{bkcenter}
\begin{bktable}{colspec={Q[l,m,wd=30mm] *{5}{Q[c,m,wd=16mm]}}, rowsep=1.5pt}
時間 $t$（s） & 0 & 1 & 2 & 3 & 4 \\
回った角（度） & 0° & 90° & 180° & 270° & 360° \\
回った角（rad） & 0 & $\dfrac{\pi}{2}$ & $\pi$ & $\dfrac{3\pi}{2}$ & $2\pi$ \\
\end{bktable}
\end{bkcenter}
同じ速さで回るので、回った角は時間に比例します（2-15）。1秒に $\dfrac{\pi}{2}$\,rad（90°）ずつふえるので、$t$ 秒で回る角は
\bkiviline{$\theta=\dfrac{\pi}{2}\,t$}
です。この $\dfrac{\pi}{2}$ のように、\textbf{1秒あたりに回る角}を\textbf{角速度}（かくそくど）といい、$\omega$（オメガ）で表します。\bkidx{かくそくど}{角速度}単位は rad/s（ラジアン毎秒）です。角速度 $\omega$ で回る点が $t$ 秒で回る角は、
\bkiviframe[60mm]{\large $\theta=\omega t$}
「速さ×時間＝道のり」と同じ形です。角速度は「角の速さ」です。

\bkivihead{周期 $T$ と振動数 $f$}
1周にかかる時間を、回るときの\textbf{周期}といい、$T$ で表します（③の周期と同じことばです。時間でいうと、同じ動きをくり返す1回分の時間です）。1周は $2\pi$\,rad で、1秒に $\omega$\,rad ずつ回るので、$\omega T=2\pi$。両辺を $\omega$ でわって（2-14）、
\bkiviline{$T=\dfrac{2\pi}{\omega}$}
回転台なら $\omega=\dfrac{\pi}{2}$ なので、$T=2\pi\div\dfrac{\pi}{2}=2\pi\times\dfrac{2}{\pi}=4$ で、4秒。合っています。

1秒あたりに回る回数を\textbf{振動数}（しんどうすう）といい、$f$ で表します。単位はHz（ヘルツ）で、「1秒に何回」という意味です。1周に $T$ 秒かかるので、1秒では $\dfrac{1}{T}$ 回です。また、1回転は $2\pi$\,rad なので、1秒に $f$ 回転すると、$2\pi\times f$\,rad 回ります。
\bkiviline{$f=\dfrac{1}{T}$\qquad $\omega=2\pi f$}
回転台なら、$f=\dfrac{1}{4}=0.25$ で\qty{0.25}{Hz}。$\omega=2\pi\times0.25=\dfrac{\pi}{2}$ で、合っています。

\bkivihead{回る点の位置と、影のグラフ}
半径 $r$ の円の上を、角速度 $\omega$ で回る点は、$t$ 秒後に角 $\omega t$ の所にいます（0秒で $x$ 軸の上から出発したとき）。②の「半径 $r$ の円のとき」の $\theta$ に $\omega t$ を入れると、
\bkiviframe[90mm]{\large $(r\cos\omega t,\ r\sin\omega t)$}
影の高さ $y=r\sin\omega t$ を、横軸を時間 $t$ にしてグラフにすると、振幅 $r$、周期 $T$ の波になります。
\begin{bkfigure}
\begin{tikzpicture}
\begin{axis}[bkplot, width=110mm, height=50mm, xmin=0, xmax=9.2, ymin=-1.35, ymax=1.45, xtick={1,2,3,4,5,6,7,8}, xticklabels={}, ytick={-1,1}, yticklabels={$-r$,$r$}, xlabel={$t$（s）}, ylabel={$y$}, tick label style={font=\sffamily\footnotesize}]
  \foreach \v in {1,...,8} {\edef\tmp{\noexpand\node[font=\noexpand\sffamily\noexpand\footnotesize, anchor=north, inner sep=1pt] at (axis cs:\v,-1.37) {\v};}\tmp}
  \addplot[bkcurve, domain=0:8.6, samples=160] {sin(90*x)};
  \draw[bkOrange!85!black, line width=0.9pt, {Stealth[length=4pt]}-{Stealth[length=4pt]}] (axis cs:1,1.18) -- (axis cs:5,1.18);
  \node[font=\sffamily\footnotesize, anchor=south, inner sep=1pt] at (axis cs:3,1.2) {周期 $T=4$\,s};
  \draw[bkGray, line width=0.5pt, dash pattern=on 2pt off 1.5pt] (axis cs:1,1) -- (axis cs:1,1.12) (axis cs:5,1) -- (axis cs:5,1.12);
\end{axis}
\end{tikzpicture}
\bkfigcaption{回転台（$T=\qty{4}{s}$、$\omega=\frac{\pi}{2}$\,rad/s）のふちの点の高さ $y=r\sin\omega t$。山から次の山まで4秒。}
\end{bkfigure}

\bkivihead{回る点の速さ $v=r\omega$}
回る点は、1秒に $\omega$\,rad だけ回ります。半径 $r$ の円で、角 $\omega$ の弧の長さは、4-2の「弧の長さ＝半径×角（$s=r\theta$）」から $r\omega$ です。だから、1秒に進む道のり、つまり速さは
\bkiviframe[60mm]{\large $v=r\omega$}
別の考え方でも確かめます。1周の道のりは円周の $2\pi r$ で、それを $T$ 秒で進むので、速さは $\dfrac{2\pi r}{T}=r\times\dfrac{2\pi}{T}=r\omega$。同じになりました。

ここで、角をラジアンではかる理由がわかります。ラジアンなら、「弧＝半径×角」がそのまま使えるので、$v=r\omega$ ときれいな形になります。度ではかると、$v=r\omega\times\dfrac{\pi}{180}$ のように、余分な数がついてしまいます。

\bkivihead{\bkhatenaref の答え}
\noindent\begin{minipage}[c]{0.6\linewidth}
\raggedright
半径\qty{50}{m}、1周20分の観覧車です。1周 $360°$ を20分で回るので、1分に $360°\div20=18°$ 回ります。いちばん下から乗って $t$ 分後には、いちばん下から $\theta=18°\times t$ だけ回っています。

いちばん下は、中心から見て $-90°$ の位置です。だから $\theta$ 回ったゴンドラの動径は、$\theta-90°$ の所にあります。中心からの高さは $50\sin(\theta-90°)$ で、中心はいちばん下より\qty{50}{m}高いので、いちばん下からの高さは
\[ y=50\sin(\theta-90°)+50 \]
⑦の言い方では、$y=50\sin\theta$ を右へ90°、上へ50動かしたグラフです。
\end{minipage}\par\vspace{4pt}\hfill
\begin{minipage}[c]{0.36\linewidth}
\centering
\begin{tikzpicture}[x=0.88mm, y=0.88mm]
  \draw[bkGray, line width=0.8pt] (-24,0) -- (24,0);
  \draw[bkInk!50, line width=0.8pt] (0,20) circle (20);
  \draw[bkGray, dashed, line width=0.5pt] (0,20) -- (0,0);
  \def\th{126}
  \draw[bkInk, line width=0.8pt] (0,20) -- ($(0,20)+({\th-90}:20)$);
  \draw[bkOrange, line width=1.4pt] (0,20) -- (0,{20+20*sin(\th-90)});
  \draw[bkGray, line width=0.5pt] (0,{20+20*sin(\th-90)}) -- ($(0,20)+({\th-90}:20)$);
  \fill[bkRoute2] ($(0,20)+({\th-90}:20)$) ++(-1.8,-2.6) rectangle ++(3.6,3);
  \fill[bkInk] (0,20) circle (0.8);
  \draw[bkInk, line width=0.5pt] (0,14) arc (-90:{\th-90}:6);
  \node[font=\sffamily\footnotesize] at ($(0,20)+({(\th-180)/2}:8.8)$) {$\theta$};
  \draw[bkGray, line width=0.5pt, <->] (-22,0) -- (-22,20);
  \node[font=\sffamily\footnotesize, anchor=east] at (-22.8,10) {50};
  \node[font=\sffamily\footnotesize, anchor=north, inner sep=1pt] at (0,-0.6) {いちばん下};
\end{tikzpicture}
\end{minipage}

\noindent 7分後は $\theta=18°\times7=126°$ で、$\theta-90°=36°$。$\sin36°=0.588$ なので、
\bkiviline{$y=50\times0.588+50=29.4+50=79.4$}
で、約\qty{79}{m}です。105分後は、$105\div20=5$ あまり5 なので、5周と5分。5周すると、いちばん下にもどるので、あまりの5分だけ考えます。$\theta=18°\times5=90°$ で、$y=50\sin0°+50=\qty{50}{m}$。中心と同じ高さです。

この式は、⑤を使うと、もっと短く書けます。⑤の $\cos(\theta+90°)=-\sin\theta$ で、$\theta$ のかわりに $\theta-90°$ を入れると、$\cos\theta=-\sin(\theta-90°)$、つまり $\sin(\theta-90°)=-\cos\theta$ です。だから
\bkiviline{$y=50-50\cos\theta$}
とも書けます。7分後なら $50-50\cos126°=50-50\times(-0.588)=79.4$ で、同じ答えです。5-15では、この形を使います。

観覧車の速さも求めます。1周の時間は $20\times\qty{60}{s/min}=\qty{1200}{s}$ です。だから、角速度は $\omega=\dfrac{2\pi}{1200}\approx\qty{0.00524}{rad/s}$ で、速さは $v=r\omega$ から、$50\times0.00524\approx\qty{0.26}{m/s}$ です。人が歩く速さ（約\qty{1.3}{m/s}）の $\dfrac{1}{5}$ ほどの、ゆっくりした速さです。



\begin{motto}[4]
\textbf{円をつぶすと、だ円}\par
単位円の上の点 $(\cos t,\ \sin t)$ を考えます。その横を $a$ 倍、たてを $b$ 倍にした点は、$(a\cos t,\ b\sin t)$ です。$t$ を0から $2\pi$ まで変えると、この点は、円を横に $a$ 倍、たてに $b$ 倍にのばした（つぶした）形をえがきます。この形を\textbf{だ円}（だえん）といいます。$\cos t=\dfrac{x}{a}$、$\sin t=\dfrac{y}{b}$ を $\sin^2t+\cos^2t=1$ に入れると、だ円の式 $\Bigl(\dfrac{x}{a}\Bigr)^2+\Bigl(\dfrac{y}{b}\Bigr)^2=1$ が出てきます。

このように、$x$ と $y$ を、別の変数 $t$ の式で表すやり方を、\textbf{媒介変数}（ばいかいへんすう）表示といいます。$t$ が、$x$ と $y$ のなかだちをしています。$x$ と $y$ でちがう角速度にした $(\sin t,\ \sin2t)$ などの形は、\textbf{リサージュ図形}といいます。2つのゆれを組み合わせたときにできる形で、電気の波をくらべる機械（オシロスコープ）の画面で見られます。

\begin{bkcenter}
\begin{tikzpicture}[x=1mm, y=1mm]
  \draw[bkGray, line width=0.5pt] (-24,0) -- (24,0) (0,-15) -- (0,15);
  \draw[bkInk!40, line width=0.8pt, dash pattern=on 2pt off 1.5pt] (0,0) circle (12);
  \draw[bkBlue, line width=1.5pt, domain=0:360, samples=120, smooth, variable=\t] plot ({20*cos(\t)},{12*sin(\t)});
  \node[font=\sffamily\footnotesize, anchor=north] at (0,-16) {$(2\cos t,\ 1.2\sin t)$ のだ円};
  \begin{scope}[xshift=62mm]
  \draw[bkGray, line width=0.5pt] (-18,0) -- (18,0) (0,-15) -- (0,15);
  \draw[bkOrange, line width=1.5pt, domain=0:360, samples=200, smooth, variable=\t] plot ({14*sin(\t)},{12*sin(2*\t)});
  \node[font=\sffamily\footnotesize, anchor=north] at (0,-16) {$(\sin t,\ \sin2t)$ のリサージュ図形};
  \end{scope}
\end{tikzpicture}
\end{bkcenter}
地球が太陽のまわりを回る道すじも、だ円です。ただし、太陽はだ円の中心ではなく、少しずれた所にあります（4-15・物理シリーズでくわしく）。
\end{motto}

\begin{taisetsu}
\def\tsl#1{\makebox[5.4\zw][l]{\sffamily\bfseries\small\color{\bkTT}#1}}
\tsl{単位円}単位円の上の点の座標は $(\cos\theta,\ \sin\theta)$\\[3pt]
\tsl{極座標}$(r,\ \theta)$ の点は $x=r\cos\theta$、$y=r\sin\theta$（どの向きでも）\\[3pt]
\tsl{一般角}$\theta+360°\times n$ は、$\theta$ と同じ位置。周期は360°（$2\pi$）\\[3pt]
\tsl{三平方}$\sin^2\theta+\cos^2\theta=1$\\[3pt]
\tsl{90°回す}$(x,\ y)\to(-y,\ x)$\\[3pt]
\tsl{}$\cos(\theta+90°)=-\sin\theta$、$\sin(\theta+90°)=\cos\theta$\\[3pt]
\tsl{回る点}$\theta=\omega t$、$(r\cos\omega t,\ r\sin\omega t)$、$T=\dfrac{2\pi}{\omega}$、$f=\dfrac{1}{T}$、$v=r\omega$
\end{taisetsu}

\begin{kotoba}
\small
\bkword{一般角}{360°より大きい角や、マイナスの角までふくめて考えた角。左回りが正。}
\bkword{動径}{始線（$x$ 軸の正の部分）から回っていく半直線OP。}
\bkword{単位円}{原点を中心とする半径1の円。}
\bkword{象限}{$x$ 軸と $y$ 軸で分けた平面の4つの区画。右上が第1象限で、左回りに第2・第3・第4。}
\bkword{三角関数}{$\sin\theta$、$\cos\theta$、$\tan\theta$ のこと。角 $\theta$ の関数。}
\bkword{周期}{同じ変わり方をくり返すときの、1回分のはば（角や時間）。}
\bkword{振幅}{波の中心の線から山までの高さ。}
\bkword{角速度}{1秒あたりに回る角 $\omega$（rad/s）。}
\bkword{振動数}{1秒あたりにくり返す回数 $f$（Hz）。$f=\dfrac{1}{T}$。}
\end{kotoba}

\bkkigocards{\bkkigocard{$\omega$}{オメガ}{角速度。1秒に回る角（rad/s）}{angular velocity}{4-8}}
{\bkkigocard{$\sin^2\theta$}{サイン・にじょう・シータ}{$(\sin\theta)^2$ のこと。$\sin(\theta^2)$ ではない}{sine squared}{4-8}}



\begin{column}{時計の針と地球の角速度}
いろいろなものの角速度を、$\omega=\dfrac{2\pi}{T}$ で計算してみます。
\begin{bkcenter}
\begin{bktable}{colspec={Q[l,m,wd=26mm] Q[c,m,wd=30mm] Q[c,m,wd=40mm]}, rowsep=1.5pt}
回るもの & 周期 $T$ & 角速度 $\omega$（rad/s） \\
時計の秒針 & 1分（\qty{60}{s}） & $\dfrac{2\pi}{60}\approx0.105$ \\
時計の長針 & 1時間（\qty{3600}{s}） & $\dfrac{2\pi}{3600}\approx0.00175$ \\
地球の自転 & 1日（\qty{86400}{s}） & $\dfrac{2\pi}{\num{86400}}\approx0.0000727$ \\
\end{bktable}
\end{bkcenter}
地球の角速度はとても小さいのですが、半径が大きいので、赤道の上の地面は $v=r\omega$ で、$\num{6400000}\times0.0000727\approx\qty{465}{m/s}$ もの速さで動いています（地球の半径は約\qty{6400}{km}）。わたしたちは、地面といっしょに動いているので、その速さを感じないのです。（地球がまわりの星に対して1周する時間は、くわしくは23時間56分ほどです。）
\end{column}


\begin{column}{自転車のギアと角速度}
\noindent\begin{minipage}[c]{0.62\linewidth}
\raggedright
自転車では、ペダルの前ギア（半径\qty{0.10}{m}）と後ろのギア（半径\qty{0.04}{m}）がチェーンでつながっています。チェーンは1本なので、2つのギアのふちは同じ速さ $v$ で動きます。$v=r\omega$ なので、
\bkiviline{$0.10\times\omega_{\text{前}}=0.04\times\omega_{\text{後}}$}
後ろのギアの角速度は、前の $0.10\div0.04=2.5$ 倍です。ペダルを1秒に1回転こぐと、後ろの車輪は1秒に2.5回転し、半径\qty{0.33}{m}のタイヤなら、$0.33\times2.5\times2\pi\approx\qty{5.2}{m/s}$（時速約\qty{19}{km}）で進みます。
\end{minipage}\hfill
\begin{minipage}[c]{0.34\linewidth}
\centering
\begin{tikzpicture}[x=0.82mm, y=0.82mm]
  \draw[bkInk!70, line width=0.9pt] (0,0) circle (10);
  \draw[bkInk!70, line width=0.9pt] (34,0) circle (4);
  \draw[bkGray, line width=1.4pt] (0,10) -- (34,4) (0,-10) -- (34,-4);
  \fill[bkInk] (0,0) circle (0.8) (34,0) circle (0.6);
  \draw[bkOrange, line width=1pt, -{Stealth[length=3.5pt]}] (12,9.4) -- (22,7.6);
  \node[font=\sffamily\footnotesize, anchor=north, inner sep=1pt] at (0,-11.5) {前ギア};
  \node[font=\sffamily\footnotesize, anchor=north, inner sep=1pt] at (34,-5.5) {後ろギア};
  \node[font=\sffamily\footnotesize, text=bkOrange!80!black, anchor=south, inner sep=1pt] at (17,9.8) {$v$};
\end{tikzpicture}
\end{minipage}
\end{column}

%% ---------------- 例題 ----------------
\bksection{reidai}{例題でたしかめよう}

\begin{reidai}[一般角を0°〜360°に直す]
次の角の動径は、0°から360°までのどの角と同じ位置ですか。何周と何度かも答えましょう。\par
（1）750°\quad（2）$-120°$\quad（3）$\dfrac{9\pi}{4}$\,rad\par
\kangaekata 360°（ラジアンでは $2\pi$）ごとに同じ位置にもどるので、360°でわったあまりを見ます。マイナスの角は、360°をたして直します。\par
\shiki\hangindent=4.3\zw\hangafter=1 （1）$750\div360=2$ あまり30\par
\bkivib 確かめ：$360\times2+30=720+30=750$\par
\bkiviq （2）$-120°+360°=240°$\par
\bkivib 右回りに120°は、左回りに240°と同じ位置。\par
\bkiviq （3）$\dfrac{9\pi}{4}=\dfrac{8\pi}{4}+\dfrac{\pi}{4}=2\pi+\dfrac{\pi}{4}$\par
\bkivib $2\pi$ が1周なので、1周と $\dfrac{\pi}{4}$。$\dfrac{\pi}{4}$\,rad は45°（$\pi$\,rad が180°だから）。\par
\kotae\hangindent=4.3\zw\hangafter=1 （1）30°（2周と30°）\quad（2）240°\quad（3）45°（1周と $\dfrac{\pi}{4}$）
\end{reidai}

\begin{jibun}
次の角の動径は、0°から360°までのどの角と同じ位置ですか。何周と何度かも答えましょう。\par
（1）1000°\quad（2）$-45°$\quad（3）$\dfrac{7\pi}{3}$\,rad
\bkjans{（1）\nobreak $1000\div360=2$ あまり280。280°（2周と280°）。（2）\nobreak $-45°+360°=315°$。（3）\nobreak $\dfrac{7\pi}{3}=2\pi+\dfrac{\pi}{3}$ で、60°（1周と $\dfrac{\pi}{3}$）。}
\end{jibun}

\begin{reidai}[単位円で $\sin$・$\cos$ を求める]
次の角の $\cos\theta$ と $\sin\theta$ を求めましょう。\quad（1）210°\quad（2）315°\quad（3）480°\par
（4）極座標が $(4,\ 210°)$ の点と $(6,\ 315°)$ の点の、$(x,\ y)$ を求めましょう。\par
\kangaekata （4）は $x=r\cos\theta$、$y=r\sin\theta$（②）に（1）（2）の値を入れます。（1）〜（3）は、②の①〜④の順。動径をかき、$x$ 軸とのあいだのとがった角の三角形で横とたてを出し、場所で符号をつけます。\par
\shiki\hangindent=4.3\zw\hangafter=1 （1）210°は、$x$ 軸の左の部分から30°下。第3象限（左下）。\par
\bkivib 30°の三角形で、横 $\dfrac{\sqrt3}{2}$、たて $\dfrac{1}{2}$。左にも下にもいるので、どちらも負。\par
\bkiviq （2）315°は、$x$ 軸の右の部分から45°下。第4象限（右下）。\par
\bkivib 45°の三角形で、横もたても $\dfrac{\sqrt2}{2}$。右にいるので $\cos$ は正、下にいるので $\sin$ は負。\par
\bkiviq （3）$480°-360°=120°$。120°は、$x$ 軸の左の部分から60°上。第2象限（左上）。\par
\bkivib 60°の三角形で、横 $\dfrac{1}{2}$、たて $\dfrac{\sqrt3}{2}$。左にいるので $\cos$ は負。\par
\bkiviq （4）$(4,\ 210°)$：$x=4\times\Bigl(-\dfrac{\sqrt3}{2}\Bigr)=-2\sqrt3$、\enspace $y=4\times\Bigl(-\dfrac{1}{2}\Bigr)=-2$\par
\bkivib $(6,\ 315°)$：$x=6\times\dfrac{\sqrt2}{2}=3\sqrt2$、\enspace $y=6\times\Bigl(-\dfrac{\sqrt2}{2}\Bigr)=-3\sqrt2$\par
\kotae\hangindent=4.3\zw\hangafter=1 （1）$\cos210°=-\dfrac{\sqrt3}{2}$、$\sin210°=-\dfrac{1}{2}$\par
\bkiviq （2）$\cos315°=\dfrac{\sqrt2}{2}$、$\sin315°=-\dfrac{\sqrt2}{2}$\par
\bkiviq （3）$\cos480°=-\dfrac{1}{2}$、$\sin480°=\dfrac{\sqrt3}{2}$\par
\bkiviq （4）$(-2\sqrt3,\ -2)$、$(3\sqrt2,\ -3\sqrt2)$
\end{reidai}

\begin{jibun}
次の角の $\cos\theta$ と $\sin\theta$ を求めましょう。\quad（1）240°\quad（2）$-60°$\quad（3）390°\quad（4）極座標が $(6,\ 240°)$ の点と $(8,\ -60°)$ の点の、$(x,\ y)$ を求めましょう。
\bkjans{（1）\nobreak $\cos240°=-\dfrac{1}{2}$、$\sin240°=-\dfrac{\sqrt3}{2}$。（2）\nobreak $-60°$ は300°と同じ位置。$\cos(-60°)=\dfrac{1}{2}$、$\sin(-60°)=-\dfrac{\sqrt3}{2}$。（3）\nobreak 30°と同じ位置。$\cos390°=\dfrac{\sqrt3}{2}$、$\sin390°=\dfrac{1}{2}$。（4）\nobreak $(6,\ 240°)$ は $x=6\cos240°=-3$、$y=6\sin240°=-3\sqrt3$。$(8,\ -60°)$ は $x=8\cos(-60°)=4$、$y=8\sin(-60°)=-4\sqrt3$。}
\end{jibun}

\begin{bkfigure}
\begin{tikzpicture}[x=1mm, y=1mm]
  % 1 を 3.6 mm。P(4, 210°)＝(−3.46, −2)、Q(6, 315°)＝(4.24, −4.24)
  \bkiviaxes{30}{12}
  \draw[bkInk, line width=1pt] (0,0) -- (210:14.4);
  \draw[bkInk, line width=1pt] (0,0) -- (315:21.6);
  \draw[bkBlue, line width=1.3pt] (0,0) -- ({14.4*cos(210)},0);
  \draw[bkOrange, line width=1.3pt] ({14.4*cos(210)},0) -- (210:14.4);
  \draw[bkBlue, line width=1.3pt] (0,0) -- ({21.6*cos(315)},0);
  \draw[bkOrange, line width=1.3pt] ({21.6*cos(315)},0) -- (315:21.6);
  \fill[bkInk] (210:14.4) circle (1) (315:21.6) circle (1);
  \draw[bkInk, line width=0.5pt, -{Stealth[length=2.5pt]}] (4,0) arc (0:210:4);
  \draw[bkInk, line width=0.5pt, -{Stealth[length=2.5pt]}] (6.5,0) arc (0:315:6.5);
  \node[font=\sffamily\footnotesize, anchor=east, align=right] at ($(210:14.4)+(-1.5,0)$) {$(4,\ 210°)$\\$(-2\sqrt3,\ -2)$};
  \node[font=\sffamily\footnotesize, anchor=west, align=left] at ($(315:21.6)+(1.5,0)$) {$(6,\ 315°)$\\$(3\sqrt2,\ -3\sqrt2)$};
\end{tikzpicture}
\bkfigcaption{例題2（4）の2つの点。青い横の線が $x=r\cos\theta$、オレンジのたての線が $y=r\sin\theta$。点が $y$ 軸より左なら $x$ が負、$x$ 軸より下なら $y$ が負になる。}
\end{bkfigure}

\begin{reidai}[$\sin^2\theta+\cos^2\theta=1$ を使う]
$\sin\theta=0.6$ で、$\theta$ は第2象限の角です。$\cos\theta$ を求めましょう。\par
\kangaekata $\sin^2\theta+\cos^2\theta=1$ から $\cos^2\theta$ を求め、2つの候補のうち、象限に合う符号のほうを選びます。\par
\shiki\hangindent=2.6\zw\hangafter=1 $\cos^2\theta=1-\sin^2\theta=1-0.6^2=1-0.36=0.64$\par
\bkivia 2乗して0.64になる数は、0.8と $-0.8$\par
\bkivia 第2象限（左上）では、Pは原点より左にいるので、$\cos\theta$ は負\par
\bkivia 確かめ：$0.6^2+(-0.8)^2=0.36+0.64=1$\par
\kotae\hangindent=2.6\zw\hangafter=1 $\cos\theta=-0.8$
\end{reidai}



\begin{reidai}[原点のまわりに90°回す]
点 $(4,\ 2)$ を、原点のまわりに左回りに90°回した点Aを求めましょう。さらに90°回した点Bも求めましょう。\par
\kangaekata 90°回すと $(x,\ y)\to(-y,\ x)$。$x$ に4、$y$ に2を入れます。\par
\shiki\hangindent=2.6\zw\hangafter=1 A：$(-y,\ x)=(-2,\ 4)$\par
\bkivia B：Aの $x$ が $-2$、$y$ が4なので、$(-y,\ x)=(-4,\ -2)$\par
\bkivia 確かめ：Bは180°回した点なので、もとの点の横も高さも反対の $(-4,\ -2)$。合っている。\par
\kotae\hangindent=2.6\zw\hangafter=1 A $(-2,\ 4)$、B $(-4,\ -2)$
\end{reidai}

\noindent\begin{minipage}[t]{0.485\linewidth}\vspace{0pt}
\begin{jibun}
$\cos\theta=0.6$ で、$\theta$ は第4象限の角です。$\sin\theta$ を求めましょう。
\bkjans{$\sin^2\theta=1-0.36=0.64$。第4象限（右下）では $\sin\theta$ は負なので、$\sin\theta=-0.8$。}
\end{jibun}
\end{minipage}\hfill
\begin{minipage}[t]{0.485\linewidth}\vspace{0pt}
\begin{jibun}
点 $(3,\ -1)$ を、原点のまわりに左回りに90°回した点Aと、さらに90°回した点Bを求めましょう。
\bkjans{A：$(-y,\ x)=(1,\ 3)$。B：$(-3,\ 1)$。}
\end{jibun}
\end{minipage}

\begin{reidai}[グラフの振幅・周期・ずれ]
（1）$y=3\sin(\theta-30°)$ の振幅、周期と、$y=\sin\theta$ のグラフをどう動かしたかを答えましょう。また、$\theta=120°$ のときの $y$ の値を求めましょう。\par
（2）$y=\sin2\theta$ の周期を求めましょう。\par
\kangaekata ⑦の表を使います。前の数が振幅、かっこの中の「$-30°$」は右へ30°。\par
\shiki\hangindent=4.3\zw\hangafter=1 （1）振幅3。周期は $y=\sin\theta$ と同じ360°。\par
\bkivib たてに3倍にして、右へ30°動かしたグラフ。\par
\bkivib $\theta=120°$ のとき $y=3\sin(120°-30°)=3\sin90°=3\times1=3$\par
\bkiviq （2）$2\theta$ が360°になるのは $\theta=180°$ のとき。$360°\div2=180°$\par
\kotae\hangindent=4.3\zw\hangafter=1 （1）振幅3、周期360°、たてに3倍して右へ30°。$y=3$
\bkiviq （2）180°
\end{reidai}

\begin{jibun}
（1）$y=2\sin(\theta-45°)$ の振幅、周期と、$y=\sin\theta$ のグラフをどう動かしたかを答えましょう。また、$\theta=135°$ のときの $y$ の値を求めましょう。\par
（2）$y=\sin3\theta$ の周期を求めましょう。
\bkjans{（1）\nobreak 振幅2、周期360°。たてに2倍して右へ45°。$y=2\sin90°=2$。（2）\nobreak $360°\div3=120°$。}
\end{jibun}

\begin{reidai}[$\sin\theta=k$、$\cos\theta=k$ になる角]
0°から360°までで、次の式が成り立つ $\theta$ を求めましょう。\par
（1）$\sin\theta=-0.5$\quad（2）$\cos\theta=-0.5$\quad（3）$\sin\theta=0.8$（電卓を使う）\par
\kangaekata 単位円で、$\sin$ は高さの横線、$\cos$ は横の位置のたての線と交わる2点を探します。\par
\shiki\hangindent=4.3\zw\hangafter=1 （1）高さ $-0.5$ の横線と交わる点。どちらも $x$ 軸の下で、30°の三角形ができる。\par
\bkivib 左下は $180°+30°=210°$、右下は $360°-30°=330°$\par
\bkiviq （2）$x=-0.5$ のたての線と交わる点。上は120°（$\cos120°=-0.5$）。\par
\bkivib 下は、$x$ 軸で裏返した $360°-120°=240°$\par
\bkiviq （3）電卓（DEG）で $\sin^{-1}0.8=53.13\cdots$ なので、約53.1°。\par
\bkivib もう1つは $180°-53.1°=126.9°$\par
\kotae\hangindent=4.3\zw\hangafter=1 （1）210°、330°\quad（2）120°、240°\quad（3）約53.1°、約126.9°
\end{reidai}

\begin{jibun}
0°から360°までで、次の式が成り立つ $\theta$ を求めましょう。\par
（1）$\sin\theta=\dfrac{\sqrt3}{2}$\quad（2）$\cos\theta=0.5$\quad（3）$\sin\theta=0.3$（電卓を使う）
\bkjans{（1）\nobreak 60°、120°。（2）\nobreak 60°、300°。（3）\nobreak $\sin^{-1}0.3=17.46\cdots$ で約17.5°、もう1つは約162.5°。}
\end{jibun}

\noindent\begin{minipage}[c]{0.36\linewidth}
\centering
\begin{tikzpicture}[x=1mm, y=1mm]
  \def\R{16}
  \bkiviaxes{21}{21}
  \draw[bkInk!45, line width=0.8pt] (0,0) circle (\R);
  \draw[bkOrange, line width=1pt] (-8,-20) -- (-8,20);
  \foreach \a in {120,240} {\draw[bkInk, line width=0.8pt] (0,0) -- (\a:\R); \fill[bkInk] (\a:\R) circle (1.1);}
  \node[font=\sffamily\footnotesize, anchor=south east, inner sep=1pt] at ($(120:\R)+(-0.5,0.5)$) {120°};
  \node[font=\sffamily\footnotesize, anchor=north east, inner sep=1pt] at ($(240:\R)+(-0.5,-0.5)$) {240°};
  \node[font=\sffamily\footnotesize, text=bkOrange!80!black, anchor=north west, inner sep=1pt] at (-7.2,-17.5) {$x=-0.5$};
\end{tikzpicture}
\end{minipage}\hfill
\begin{minipage}[c]{0.6\linewidth}
\raggedright
\textbf{例題6の（2）を図で見る}\par
$\cos\theta$ は横の位置なので、$\cos\theta=-0.5$ になる点は、たての線 $x=-0.5$（オレンジ）と単位円が交わる所です。上の点が120°、下の点は、上の点を $x$ 軸で裏返した位置で、$-120°$、つまり240°です。電卓の $\cos^{-1}(-0.5)$ は120°だけを教えてくれます。
\end{minipage}

\begin{reidai}[角速度・周期・振動数・速さ]
扇風機の羽根が、1分に1200回転しています。羽根の先は、中心から\qty{0.15}{m}です。振動数 $f$、周期 $T$、角速度 $\omega$、羽根の先の速さ $v$ を求めましょう。\par
\kangaekata 1分を秒に直して、1秒あたりの回転数（振動数）を出します。$T=\dfrac{1}{f}$、$\omega=2\pi f$、$v=r\omega$。\par
\shiki\hangindent=2.6\zw\hangafter=1 $f=\dfrac{1200\text{ 回}}{\qty{60}{s}}=20$ で\qty{20}{Hz}\par
\bkivia $T=\dfrac{1}{20}=0.05$ で\qty{0.05}{s}\par
\bkivia $\omega=2\pi\times20=40\pi\approx125.7$ で約\qty{126}{rad/s}\par
\bkivia $v=r\omega=0.15\times125.7\approx18.8$ で約\qty{18.8}{m/s}\par
\bkivia 確かめ：$\omega=\dfrac{2\pi}{T}=\dfrac{2\pi}{0.05}=40\pi$。同じ。\par
\kotae\hangindent=2.6\zw\hangafter=1 $f=\qty{20}{Hz}$、$T=\qty{0.05}{s}$、$\omega=40\pi\approx\qty{126}{rad/s}$、$v\approx\qty{18.8}{m/s}$
\end{reidai}
羽根の先の速さ\qty{18.8}{m/s}は、時速になおすと約\qty{68}{km/h}です（\qty{1}{m/s}は\qty{3.6}{km/h}）。回っている扇風機に指を入れてはいけないわけが、数でわかります。



\begin{jibun}
洗濯機のドラムが、脱水のとき1分に600回転しています。ドラムの半径は\qty{0.25}{m}です。$f$、$T$、$\omega$、ドラムのふちの速さ $v$ を求めましょう。
\bkjans{$f=600\div60=\qty{10}{Hz}$、$T=\qty{0.1}{s}$、$\omega=2\pi\times10=20\pi\approx\qty{62.8}{rad/s}$、$v=0.25\times62.8\approx\qty{15.7}{m/s}$。}
\end{jibun}

\begin{reidai}[観覧車の高さ]
半径\qty{30}{m}、1周12分の観覧車に、いちばん下から乗ります。いちばん下からのゴンドラの高さは、$\theta$ 回ったとき $y=30-30\cos\theta$（m）です。\par
（1）1分に何度回りますか。\quad（2）4分後の高さは。\quad（3）50分後の高さは。\par
\kangaekata ⑨の \bkhatenaref の答えと同じ式です（半径が30）。何周もしたあとは、あまりの時間だけ考えます。\par
\shiki\hangindent=4.3\zw\hangafter=1 （1）$360°\div12=30°$\par
\bkiviq （2）$\theta=30°\times4=120°$。$\cos120°=-0.5$ なので\par
\bkivib $y=30-30\times(-0.5)=30+15=45$\par
\bkiviq （3）$50\div12=4$ あまり2 なので、4周と2分。$\theta=30°\times2=60°$\par
\bkivib $y=30-30\cos60°=30-30\times0.5=30-15=15$\par
\kotae\hangindent=4.3\zw\hangafter=1 （1）30°\quad（2）\qty{45}{m}\quad（3）\qty{15}{m}
\end{reidai}



\begin{jibun}
半径\qty{40}{m}、1周16分の観覧車に、いちばん下から乗ります。高さは $y=40-40\cos\theta$（m）です。（1）1分に何度回りますか。（2）6分後の高さは。（3）50分後の高さは（$\cos45°=0.707$、$\cos135°=-0.707$）。
\bkjans{（1）\nobreak $360°\div16=22.5°$。（2）\nobreak $\theta=135°$。$40+40\times0.707\approx\qty{68.3}{m}$。（3）\nobreak $50\div16=3$ あまり2 で、$\theta=45°$。$40-40\times0.707\approx\qty{11.7}{m}$。}
\end{jibun}


\begin{reidai}[回る点の影の高さ]
回転台のふちの点が、半径\qty{0.1}{m}の円を、周期\qty{4}{s}で左回りに回っています。0秒のとき、点は $x$ 軸の上にいます。影の高さは $y=0.1\sin\omega t$（m）です。\par
（1）$\omega$ を求めましょう。\quad（2）0.5秒後と3秒後の影の高さを求めましょう。\par
\kangaekata $\omega=\dfrac{2\pi}{T}$。$\omega t$ はラジアンなので、電卓は RAD にするか、度に直して計算します（$\pi$\,rad は180°）。\par
\shiki\hangindent=4.3\zw\hangafter=1 （1）$\omega=\dfrac{2\pi}{4}=\dfrac{\pi}{2}$\par
\bkiviq （2）0.5秒後：$\omega t=\dfrac{\pi}{2}\times0.5=\dfrac{\pi}{4}$\,rad（45°）\par
\bkivib $y=0.1\sin45°=0.1\times0.707\approx0.071$\par
\bkivib 3秒後：$\omega t=\dfrac{\pi}{2}\times3=\dfrac{3\pi}{2}$\,rad（270°）\par
\bkivib $y=0.1\sin270°=0.1\times(-1)=-0.1$\par
\kotae\hangindent=4.3\zw\hangafter=1 （1）$\dfrac{\pi}{2}$\,rad/s\quad（2）約\qty{0.071}{m}、$-\qty{0.1}{m}$（中心より\qty{0.1}{m}下）
\end{reidai}

\begin{jibun}
半径\qty{0.2}{m}、周期\qty{6}{s}で回る点の影の高さは、$y=0.2\sin\omega t$（m）です。（1）$\omega$ を求めましょう。（2）1秒後と4.5秒後の影の高さを求めましょう（$\sin60°=0.866$）。
\bkjans{（1）\nobreak $\omega=\dfrac{2\pi}{6}=\dfrac{\pi}{3}$\,rad/s。（2）\nobreak 1秒後は $\dfrac{\pi}{3}$（60°）で $0.2\times0.866\approx\qty{0.173}{m}$。4.5秒後は $\dfrac{3\pi}{2}$（270°）で $-\qty{0.2}{m}$。}
\end{jibun}

%% ---------------- まちがえやすいところ・物理工学 ----------------
\bksection{tsumazuki}{まちがえやすいところをチェック}
\noindent\begin{minipage}[t]{0.485\linewidth}\vspace{0pt}
\begin{tsumazuki}{$\cos$ の符号をわすれる}
\bkng{$\cos150°=0.866$}\bkok{$\cos150°=-0.866$}
150°の点Pは、原点より左にいます。単位円をかいて、Pの場所で符号を決めましょう。
\end{tsumazuki}
\end{minipage}\hfill
\begin{minipage}[t]{0.485\linewidth}\vspace{0pt}
\begin{tsumazuki}{答えが1つだけ}
\bkng{$\sin\theta=0.5$ なら $\theta=30°$}\bkok{$\theta=30°$、150°}
高さ0.5の点は、単位円の右と左に1つずつあります（0°〜360°のとき）。
\end{tsumazuki}
\end{minipage}
\par\vspace{2pt}
\noindent\begin{minipage}[t]{0.485\linewidth}\vspace{0pt}
\begin{tsumazuki}{ずらす向き}
\bkng{$\sin(\theta-60°)$ は左へ60°}\bkok{右へ60°}
かっこの中が0になるのは $\theta=60°$。波の始まりが右へ動きます。
\end{tsumazuki}
\end{minipage}\hfill
\begin{minipage}[t]{0.485\linewidth}\vspace{0pt}
\begin{tsumazuki}{DEG と RAD}
\bkng{DEG：$\sin1.5708=0.027$}\bkok{RAD：$\sin1.5708=1$}
$\omega t$ の角はラジアンです。答えが小さすぎたら、電卓の表示を確かめましょう。
\end{tsumazuki}
\end{minipage}

\begin{butsuri}
\bkivibox{ばねのおもりは、回る点の影と同じに動く}
\noindent\begin{minipage}[c]{0.56\linewidth}
\raggedright
回転台のふちのくぎと、ばねにつるしたおもりを、同じ周期で動かして真横から光を当てると、2つの影はぴったり同じに上下します。おもりの高さも $r\sin\omega t$ の形で変わるのです。このゆれを\textbf{単振動}といいます（6-6でくわしく）。
\end{minipage}\hfill
\begin{minipage}[c]{0.4\linewidth}
\centering
\begin{tikzpicture}[x=0.9mm, y=0.9mm]
  \draw[bkInk!55, line width=0.8pt] (0,0) circle (13);
  \fill[bkInk] (0,0) circle (0.8);
  \draw[bkInk, line width=0.8pt] (0,0) -- (40:13);
  \fill[bkOrange] (40:13) circle (1.4);
  \draw[bkInk, line width=0.6pt, -{Stealth[length=3.5pt]}] (17,-4) arc (-15:20:17);
  % ばね
  \draw[bkGray, line width=1.2pt] (24,22) -- (38,22);
  \draw[bkInk, line width=0.7pt] (31,22) -- (31,20) -- (29,19) -- (33,17.5) -- (29,16) -- (33,14.5) -- (29,13) -- (33,11.5) -- (31,10.5) -- (31,{13*sin(40)+2});
  \filldraw[fill=bkBlue!60, draw=bkBlue!70!black, line width=0.5pt] (28,{13*sin(40)-2.5}) rectangle (34,{13*sin(40)+2});
  \draw[bkOrange, line width=0.6pt, dash pattern=on 2pt off 1.5pt] (40:13) -- (27.5,{13*sin(40)});
  \draw[bkInk, line width=0.6pt, {Stealth[length=3pt]}-{Stealth[length=3pt]}] (37,-13) -- (37,13);
\end{tikzpicture}
\end{minipage}
\end{butsuri}

\begin{butsuri}
\bkivibox{エンジンのクランク}
\noindent\begin{minipage}[c]{0.56\linewidth}
\raggedright
自動車のエンジンでは、ピストンの往復の動きを、クランクという回る部品で回転に変えています。クランクのピンは半径 $r$ の円を回り、ピンの横の位置は $r\cos\theta$ です。ピストンは、ピンの横の位置（影）とほとんど同じに動きます。ずれは、ピンとピストンをつなぐ棒（連接棒）が長いほど小さくなります。
\end{minipage}\hfill
\begin{minipage}[c]{0.4\linewidth}
\centering
\begin{tikzpicture}[x=0.85mm, y=0.85mm]
  \draw[bkInk!55, line width=0.8pt] (0,0) circle (8);
  \fill[bkInk] (0,0) circle (0.8);
  \draw[bkInk, line width=1.1pt] (0,0) -- (60:8);
  \fill[bkOrange] (60:8) circle (1.2);
  \pgfmathsetmacro\px{8*cos(60)+sqrt(26^2-(8*sin(60))^2)}
  \draw[bkGray, line width=2pt] (60:8) -- (\px,0);
  \draw[bkInk!70, line width=0.8pt] (20,-6) -- (46,-6) (20,6) -- (46,6) (46,-6) -- (46,6);
  \filldraw[fill=bkBlue!50, draw=bkBlue!70!black, line width=0.5pt] ({\px-3},-5.3) rectangle ({\px+5},5.3);
  \node[font=\sffamily\footnotesize, anchor=north] at ({\px+1},-6.8) {ピストン};
  \node[font=\sffamily\footnotesize, anchor=north] at (0,-9) {クランク};
\end{tikzpicture}
\end{minipage}
\end{butsuri}

\begin{butsuri}
\bkivibox{家のコンセントの交流}
発電機の中では、磁石のあいだでコイル（巻いた電線）が回っています。コイルが角速度 $\omega$ で回ると、電圧は $V_0\sin\omega t$ の形で変わり（$V_0$ は山の高さ）、向きが半周ごとに入れかわります。これが\textbf{交流}です。

東日本では1秒に50回、西日本では60回くり返します。50\,Hz なら、周期 $T=\dfrac{1}{50}=\qty{0.02}{s}$、角速度 $\omega=2\pi\times50\approx\qty{314}{rad/s}$ です。家の「\qty{100}{V}」は実効値で、山の高さは約\qty{141}{V}です（3-10）。
\begin{bkcenter}
\begin{tikzpicture}
\begin{axis}[bkplot, width=112mm, height=29mm, xmin=0, xmax=0.046, ymin=-175, ymax=195, xtick={0.01,0.02,0.03,0.04}, xticklabels={}, ytick={-141,141}, xlabel={$t$（s）}, ylabel={電圧（V）}, tick label style={font=\sffamily\footnotesize}]
  \foreach \v in {0.01,0.02,0.03,0.04} {\edef\tmp{\noexpand\node[font=\noexpand\sffamily\noexpand\footnotesize, anchor=north, inner sep=1pt] at (axis cs:\v,-177) {\v};}\tmp}
  \addplot[bkcurve, domain=0:0.043, samples=200] {141*sin(360*50*x)};
  \draw[bkOrange!85!black, line width=0.9pt, {Stealth[length=4pt]}-{Stealth[length=4pt]}] (axis cs:0.005,160) -- (axis cs:0.025,160);
  \node[font=\sffamily\footnotesize, anchor=west, inner sep=1pt] at (axis cs:0.026,160) {周期 \qty{0.02}{s}（50\,Hz）};
\end{axis}
\end{tikzpicture}
\end{bkcenter}
\end{butsuri}



%% ---------------- 章末問題 ----------------
\bkendproblems
\bkproblemgroup{基本}{3}

\bkproblem 次の角の動径は、0°から360°までのどの角と同じ位置ですか\mbox{（\bkreidairef{1}）}。\\（1）840°\quad（2）$-210°$\quad（3）$\dfrac{17\pi}{6}$\,rad
\bkpans{（1）\nobreak 2周と120°で、120°。（2）\nobreak $-210°+360°=150°$。（3）\nobreak $\dfrac{17\pi}{6}=2\pi+\dfrac{5\pi}{6}$ で、150°。}

\bkproblem 次の角の $\cos\theta$ と $\sin\theta$ を求めましょう\mbox{（\bkreidairef{2}）}。\quad（1）225°\quad（2）$-30°$\quad（3）510°\quad（4）極座標が $(10,\ 225°)$ の点と $(4,\ 510°)$ の点の $(x,\ y)$
\bkpans{（1）\nobreak $\cos225°=-\dfrac{\sqrt2}{2}$、$\sin225°=-\dfrac{\sqrt2}{2}$。（2）\nobreak $\cos(-30°)=\dfrac{\sqrt3}{2}$、$\sin(-30°)=-\dfrac{1}{2}$。（3）\nobreak 150°と同じ位置。$\cos510°=-\dfrac{\sqrt3}{2}$、$\sin510°=\dfrac{1}{2}$。（4）\nobreak $(10,\ 225°)$ は $(-5\sqrt2,\ -5\sqrt2)$、$(4,\ 510°)$ は $(-2\sqrt3,\ 2)$。}

\bkproblem $\sin\theta=-0.8$ で、$\theta$ は第3象限の角です。$\cos\theta$ を求めましょう\mbox{（\bkreidairef{3}）}。
\bkpans{$\cos^2\theta=1-0.64=0.36$。第3象限では負なので、$\cos\theta=-0.6$。}

\bkproblem 点 $(5,\ 2)$ を、原点のまわりに左回りに90°回した点と、さらに90°回した点を求めましょう\mbox{（\bkreidairef{4}）}。
\bkpans{$(-2,\ 5)$、$(-5,\ -2)$。}

\bkproblem $y=4\sin(\theta-60°)$ の振幅、周期と、$y=\sin\theta$ のグラフをどう動かしたかを答えましょう\mbox{（\bkreidairef{5}）}。また、$\theta=150°$ のときの $y$ の値を求めましょう。
\bkpans{振幅4、周期360°。たてに4倍して右へ60°。$y=4\sin90°=4$。}

\bkproblem 0°から360°までで、次の式が成り立つ $\theta$ を求めましょう\mbox{（\bkreidairef{6}）}。\\（1）$\sin\theta=-\dfrac{\sqrt3}{2}$\quad（2）$\cos\theta=\dfrac{\sqrt2}{2}$
\bkpans{（1）\nobreak 240°、300°。（2）\nobreak 45°、315°。}

\bkproblem 自転車のタイヤが、1分に180回転しています。タイヤの半径は\qty{0.33}{m}です。$f$、$T$、$\omega$ と、タイヤのふちの速さ $v$ を求めましょう\mbox{（\bkreidairef{7}）}。
\bkpans{$f=180\div60=\qty{3}{Hz}$、$T=\dfrac{1}{3}\approx\qty{0.33}{s}$、$\omega=6\pi\approx\qty{18.8}{rad/s}$、$v=0.33\times18.85\approx\qty{6.2}{m/s}$。}

\bkproblemgroup{チャレンジ}{4}

\bkproblem 半径\qty{50}{m}、1周20分の観覧車に、いちばん下から乗ります。高さは $y=50-50\cos\theta$（m）です\mbox{（\bkreidairef{8}）}。（1）10分後の高さは。（2）65分後の高さは。（3）3分後の高さは（$\cos54°=0.588$）。
\bkpans{（1）\nobreak $\theta=180°$。$50+50=\qty{100}{m}$。（2）\nobreak $65\div20=3$ あまり5 で $\theta=90°$。\qty{50}{m}。（3）\nobreak $\theta=54°$。$50-50\times0.588=\qty{20.6}{m}$。}

\bkproblem 0°から360°までで、$\sin\theta=0.6$ が成り立つ $\theta$ を、電卓を使って求めましょう\mbox{（\bkreidairef{6}）}（小数第1位まで）。
\bkpans{$\sin^{-1}0.6=36.86\cdots$ で約36.9°、もう1つは約143.1°。}

\bkproblem パソコンのハードディスクの円板が、1分に7200回転しています。中心から\qty{0.045}{m}の所の、$f$、$T$、$\omega$、速さ $v$ を求めましょう\mbox{（\bkreidairef{7}）}。
\bkpans{$f=7200\div60=\qty{120}{Hz}$、$T=\dfrac{1}{120}\approx\qty{0.0083}{s}$、$\omega=240\pi\approx\qty{754}{rad/s}$、$v=0.045\times754\approx\qty{33.9}{m/s}$。}

\begin{furikaeri}
\begin{checklist}
\item 単位円の上の点 $(\cos\theta,\ \sin\theta)$ を求められる。\bkfurref{bk:4-8:tanien}
\item $\sin^2\theta+\cos^2\theta=1$ や90°回す式を説明できる。\bkfurref{bk:4-8:sanheihou}
\item $\omega$ と $T$ で、回る点の位置・速さを求められる。\bkfurref{bk:4-8:omega}
\end{checklist}
\end{furikaeri}

\bkfukidashi{mina}{回る点の影が、$\sin$ の波になった！　次の章では、2つの波をたし合わせてみよう。}
